\section{Reading PAGs}

\JorisTODO{This subsection is work in progress! There may be dragons (and mistakes).}

\subsubsection{Acyclic, with selection bias}

We consider the setting with input nodes and selection bias but without cycles.
Then, we can omit the qualifier ``$\sigma$-'' from the terms.
\JorisTODO{But not the qualifier ``i''!}

We first show that nodes in a MAG that have an incoming arrowhead must be non-ancestors of selection nodes,
while nodes that are part of an undirected edge must be ancestors of selection nodes, and all other nodes
may or may not be ancestor of a selection node.
This characterization is complete (absent additional information that is not encoded in the MAG).
\begin{Lem}\label{lem:representing_CADMG_of_MAG}
  Let $G$ be a CADMG with input nodes $J$ and output nodes $V^+ = V \dcup S$ and $H$ the MAG that represents $G$ given $S$.
  Define
  $$Z^+ := \{ z \in V \cup J : \exists_{y \in V \cup J}: z \tut y \in H \}$$
  and
  $$Z^- := \{ z \in V \cup J : \exists_{x \in V \cup J}: z \hus x \in H \}.$$
  Partition $V \cup J$ into subsets $A$ and $B$, such that $Z^+ \subseteq A$ and $Z^- \subseteq B$.
  Then there exists a CADMG $\tilde{G}$ with input nodes $J$ and output nodes $\tilde{V}^+ = V \dcup \tilde{S}$ such that:
  \begin{enumerate}
    \item $\tilde{G}$ is represented by $H$ given $\tilde{S}$, and
%    \item each edge $x \suh y$ for $x,y \in V\cup J$ is present in $H$ if and only if it is present in $\tilde{G}$, and
    \item for each $z \in A$: $z \in \Anc_{\tilde{G}}(\tilde{S})$, and
    \item for each $z \in B$: $z \notin \Anc_{\tilde{G}}(\tilde{S})$.
  \end{enumerate}
\end{Lem}
\begin{proof}
  We construct the CADMG $\tilde{G}$ as follows.
  It has nodes input nodes $J$ and output nodes $\tilde{V}^+$ with $\tilde{V}^+ = V \cup \tilde{S}$, where $\tilde{S} := \{s_{\{x,y\}} : x \tut y \in H\} \cup \{s_z : z \in A \sm Z^+\}$.
  We add all edges $x \suh y$ for $x,y \in V\cup J$ that are present in $H$ to $\tilde{G}$.
  For the undirected edges $x \tut y$ in $H$, we add $x \tuh s_{\{x,y\}} \hut y$ to $\tilde{G}$.
  Finally, for each $z \in A \sm Z^+$, we add the edge $z \tuh s_z$ to $\tilde{G}$.
  Note that $\tilde{G}$ is acyclic by construction.

  Since all directed edges in $H$ are also present in $\tilde{G}$, and no directed edge absent in $H$ is present in $\tilde{G}$, the ancestral relations between variables in $V \cup J$ are the same in $H$ as in $\tilde{G}$. 
  That is, for $x,y \in V \cup J$, $x \in \Anc_H(y) \iff x \in \Anc_{\tilde{G}}(y)$.

  %We now show that for $v \in Z^+ \cup B$: if $v \in \Anc_{\tilde{G}}(\tilde{S})$, then $v \in \Anc_G(S)$ (note: the converse does not hold).
  %It suffices to show this for $v \in Z^-$.
  %We now show that for $v \in Z^-$: if $v \in \Anc_{\tilde{G}}(\tilde{S})$, then $v \in \Anc_G(S)$ (note: the converse does not hold).
  We now show: for $z \in Z^-$, $z \notin \Anc_{\tilde{G}}(\tilde{S})$.
  Let $z \in Z^-$.
  Suppose $z \in \Anc_{\tilde{G}}(\tilde{S})$.
  Then there must be a directed path $z \tuh \dots \tuh w \tuh s$ in $\tilde{G}$ such that $s \in \tilde{S}$, and all other nodes on the directed path are in $V \cup J$.
  The directed path $z \tuh \dots \tuh w$ must also be present in $H$, and by construction of $\tilde{G}$, $w \in A$. 
  Since $A \cap Z^- = \emptyset$, the path must be trivial. 
  But then still $z = w \in A$, which contradicts $z \in Z^-$.
  Hence, $z \notin \Anc_{\tilde{G}}(\tilde{S})$.

  It remains to show that $H$ represents $\tilde{G}$ given $\tilde{S}$.
  For that we need to show:
  \begin{enumerate}
    \item For all $x \in V \cup J, y \in V$: $x \sus y$ is in $H$ if and only if there is an inducing path between $x$ and $y$ given $\tilde{S}$ in $\tilde{G}$;
    \item If $x \hus y$ in $H$ then $x \notin \Anc_{\tilde{G}}(\{y\} \cup \tilde{S})$;
    \item If $x \tus y$ in $H$ then $x \in \Anc_{\tilde{G}}(\{y\} \cup \tilde{S})$;
  \end{enumerate}

  First adjacency.
  Let $x \in V \cup J$, $y \in V$.
  Suppose $x$ and $y$ are not adjacent in $H$.
  We must show that in $\tilde{G}$ there is no inducing path between $x$ and $y$ given $\tilde{S}$.
  Suppose there is one.
  Pick one of minimal length.
  It must be a collider path in which every collider is in $\Anc_{\tilde{G}}(\{x,y\} \cup \tilde{S})$.

  Suppose the path contains a node of $\tilde{S}$.
  It cannot contain some $s_z$ for $z \in A \sm Z^+$, because then it had to be of the form $x \dots z \tuh s_z \hut z \dots y$ with $z$ occurring more than once.
  Hence, the path has to be of the form $x \dots a \tuh s_{\{a,b\}} \hut b \dots y$ with $a \tut b$ in $H$.
  Since only the endpoint nodes can be non-colliders, the path must be $x \tuh s_{\{x,y\}} \hut y$ with $x \tut y$ in $H$. 
  Contradiction.
  We conclude that the path can only contain nodes in $V \cup J$.

  Note that each edge on the path must then be between a pair of nodes in $V \cup J$, and hence be present in $H$ as well.
  Then each collider on the path is not in $\Anc_G(S)$, and is in $Z^-$.
  Therefore, by the result above, each collider on the path is not in $\Anc_{\tilde{G}}(\tilde{S})$.
  But this means that each collider on the path must be in $\Anc_{\tilde{G}}(\{x,y\})$, and hence in $\Anc_H(\{x,y\})$.
  So then the path is also an inducing path between $x$ and $y$ in $H$, which contradicts the maximality of $H$.

  Now suppose $x$ and $y$ are adjacent in $H$.
  We distinguish four cases corresponding to the type of the edge between $x$ and $y$ in $H$. 
  \begin{itemize}
    \item $x \tuh y$: By construction $x \tuh y$ is in $\tilde{G}$.
      Then $y \notin \Anc_G(S)$ and $y \in Z^-$, and therefore also $y \notin \Anc_{\tilde{G}}(\tilde{S})$.
      Therefore, $x \in \Anc_{\tilde{G}}(\{y\} \cup \tilde{S})$ and $y \notin \Anc_{\tilde{G}}(\{y\} \cup \tilde{S})$.
    \item $x \hut y$: This case is similar.
    \item $x \huh y$: By construction $x \huh y$ is in $\tilde{G}$.
      Then both $x, y \notin \Anc_G(S)$ and both $x, y \in Z^-$, and therefore also $x,y \notin \Anc_{\tilde{G}}(\tilde{S})$. 
      Therefore, $x \notin \Anc_{\tilde{G}}(\{y\} \cup \tilde{S})$ and $y \notin \Anc_{\tilde{G}}(\{y\} \cup \tilde{S})$.
    \item $x \tut y$: Then $x \tuh s_{\{x,y\}} \hut y$ in $\tilde{G}$.
      Therefore, $x \in \Anc_{\tilde{G}}(\{y\} \cup \tilde{S})$ and $y \in \Anc_{\tilde{G}}(\{y\} \cup \tilde{S})$.
  \end{itemize}
  In all cases, there is an inducing path between $x$ and $y$ given $\tilde{S}$ in $\tilde{G}$.
  We conclude that $H$ represents $\tilde{G}$ given $\tilde{S}$.

  By construction of $\tilde{G}$, it holds that for all $z \in A$, $z \in \Anc_{\tilde{G}}(\tilde{S})$.
  Let $z \in B$. 
  If $z \in Z^-$, then by the properties above, we already know that $z \notin \Anc_{\tilde{G}}(\tilde{S})$.
  So consider $z \in B \sm Z^-$. Then $z$ has no adjacent undirected edge in $H$, nor an adjacent edge $z \hus \dots$.
  In other words, $z$ can only have outgoing directed edges: $z \sus \dots$ in $H$ implies $z \tuh \dots$ in $H$.
  All children of $z$ in $H$ are in $Z^-$, and hence not in $\Anc_{\tilde{G}}(\tilde{S})$. 
  The only remaining possibility for $z$ to be in $\Anc_{\tilde{G}}(\tilde{S})$ is for $z$ to be a parent of $\tilde{S}$ in $\tilde{G}$.
  But that would imply $z \in A \sm Z^+$, a contradiction. Hence $z \notin \Anc_{\tilde{G}}(\tilde{S})$.
\end{proof}
The following notion was proposed originally by \cite{Zhang2008_JMLR}. 
We extend it here to our setting which allows for selection bias and input nodes.
\begin{Def}\label{def:mag_visible_edge}
  Let $H$ be a MAG with input nodes $J$ and output nodes $V$.

  Let $a \in J \cup V$ and $b \in V$.
  A directed edge $a \tuh b$ in $H$ is called \emph{visible} if $a \in J$, or if $a \in V$ and there is a node $z$ in $H$ that is not adjacent to $b$ such that:
  \begin{itemize}
    \item there is an edge $z \suh a$ in $H$, or 
    \item there is a collider path $z \suh \dots \huh a$ in $H$ that is into $a$ and every node on the path (except $z$) is a parent of $b$ in $H$. 
  \end{itemize}
  Otherwise, $a \tuh b$ is called \emph{invisible}.
\end{Def}

\begin{figure}\centering
  \begin{tikzpicture}
    \begin{scope}[xshift=0cm]
      \node[ndint] (a) at (0,0) {$a$};
      \node[ndout] (b) at (0,-1.5) {$b$};
      \draw[tuh] (a) edge (b);
    \end{scope}
    \begin{scope}[xshift=4cm]
      \node[ndout] (z) at (-1.5,0) {$z$};
      \node[ndout] (a) at (0,0) {$a$};
      \node[ndout] (b) at (0,-1.5) {$b$};
      \draw[suh] (z) edge (a);
      \draw[tuh] (a) edge (b);
      \draw[sus,red] (z) edge (b);
    \end{scope}
    \begin{scope}[xshift=12cm]
      \node[ndout] (z) at (-6.0,0) {$z$};
      \node[ndout] (c1) at (-4.5,0) {$c_1$};
      \node (dots) at (-3.0,0) {$\cdots$};
      \node[ndout] (ck) at (-1.5,0) {$c_k$};
      \node[ndout] (a) at (0,0) {$a$};
      \node[ndout] (b) at (0,-1.5) {$b$};
      \draw[suh] (z) edge (c1);
      \draw[huh] (c1) edge (dots);
      \draw[huh] (dots) edge (ck);
      \draw[huh] (ck) edge (a);
      \draw[tuh] (a) edge (b);
      \draw[sus,red,bend right] (z) edge (b);
      \draw[tuh] (c1) edge (b);
      \draw[tuh] (ck) edge (b);
    \end{scope}
  \end{tikzpicture}
  \caption{Different patterns corresponding to a visible edge $a \tuh b$.
  $z$ and $b$ must be non-adjacent (indicated in red a forbidden edge).
  Note that $\lt b,a,c_k,\dots,c_1,z \rt$ is a discriminating path for $a$ between $b$ and $z$.
  \label{fig:visible_edge}}
\end{figure}

The following lemma of \cite{Zhang2008_JMLR} stays valid in our more general setting.
The proof closely follows the original proof.
\begin{Lem}\label{lem:zhang2008_lemma_9_mag_extended}
  Let $G$ be a CADMG with input nodes $J$ and output nodes $V^+ = V \dcup S$ and $H$ a MAG that represents $G$ given $S$.
  Let $a \in V \cup J$, $b \in V$ such that $a \tuh b$ is a directed edge in $H$. 
  If $a \tuh b$ is visible, then there exists no inducing walk between $a$ and $b$ given $S$ in $G$ that is into $a$.
  In particular, $a \huh b$ is not in $G$.
\end{Lem}
\begin{proof}
  If $a \in J$, then there are no edges into $a$ in $G$ by definition, which proves the claim.
  So assume $a \in V$. We proceed by contradiction. Assume that in $G$ there exists an inducing walk between $a$ and $b$ given $S$ that is into $a$.

  Let $z$ be a node in $V \sm \{a,b\}$.
  Consider two cases:
  \begin{itemize}
    \item
      If $z \suh a$ is in $H$, then there is an inducing walk between $z$ and $a$ in $G$ given $S$ that is into $a$ (Lemma~\ref{lem:mixed_graph_represents_CDMG}).
  Concatenating this with the inducing walk between $a$ and $b$ given $S$ that is into $a$ we obtain an inducing walk between $z$ and $b$ given $S$ (note that $a$ becomes a collider that is in $\Anc_G(\{b\} \cup S)$) in $G$. 
\item
  Suppose there is a definite collider path $\pi$ in $H$ from $z$ to $a$ that is into $a$ and such that every node on the path (except $z$) is parent of $b$ in $H$. 
  For each pair of subsequent nodes $(v_i,v_{i+1})$ on $\pi$ there exists an inducing walk in $G$ between $v_i$ and $v_{i+1}$ given $S$ that is into $v_{i+1}$, and also into $v_i$ unless possibly $v_i = z$.
  Because each node other than $z$ on $\pi$ is in $\Anc_G(\{b\} \cup S)$, all these inducing walks can be concatenated into one inducing walk in $G$ between $z$ and $b$ given $S$.
  \end{itemize}

  In both cases, $z$ and $b$ are adjacent in $H$.
  Since this holds for all such $z$, the directed edge $a \to b$ in $H$ cannot be visible. 
\end{proof}

Also the next lemma of \cite{Zhang2008_JMLR} stays valid in our more general setting.
The proof closely follows the original proof.
We extend it to also cover undirected edges.
\begin{Lem}\label{lem:zhang2008_lemma_10_mag_extended}
  Let $G$ be a CADMG with input nodes $J$ and output nodes $V^+ = V \dcup S$ and $H$ the MAG that represents $G$ given $S$.
  %Let $Z \subseteq V \cup J$ such that each of which has no incoming arrowhead in $H$ (one can always take $Z = \emptyset$).
  %Let $i \tuh j$ be a directed edge in $H$.
  If $i \tuh j$ is an invisible directed edge in $H$, or if $i \tut j$ is an undirected edge in $H$,
  then there exists a CADMG $\tilde{G}$ with input nodes $J$ and output nodes $\tilde{V}^+ = V \dcup \tilde{S}$ such that:
  \begin{enumerate}
    \item $\tilde{G}$ is also represented by $H$ given $\tilde{S}$, and
    \item $i \huh j$ is in $\tilde{G}$.
  \end{enumerate}
\end{Lem}
\begin{proof}
  We construct CADMG $\tilde{G}$ as follows.
  It has nodes $\tilde{V}^+ \cup J$ with $\tilde{V}^+ = V \cup \tilde{S}$, where $\tilde{S} := \{s_{\{x,y\}} : x \tut y \in H\}$.
  We add all edges $x \suh y$ for $x,y \in V\cup J$ that are present in $H$ to $\tilde{G}$.
  In addition, we add the bidirected edge $i \huh j$ to $\tilde{G}$.
  For the undirected edges $x \tut y$ in $H$, we add $x \tuh s_{\{x,y\}} \hut y$ to $\tilde{G}$.
  Note that $\tilde{G}$ is acyclic by construction.

  Since all directed edges in $H$ are also present in $\tilde{G}$, and no directed edge absent in $H$ is present in $\tilde{G}$, the ancestral relations between variables in $V \cup J$ are the same in $H$ as in $\tilde{G}$. 
  That is, for $x,y \in V \cup J$, $x \in \Anc_H(y) \iff x \in \Anc_{\tilde{G}}(y)$.
  
  We now show that for $v \in V \cup J$: if $v \in \Anc_{\tilde{G}}(\tilde{S})$, then $v \in \Anc_G(S)$ (note: the converse does not hold).
  Indeed, if $v \in \Anc_{\tilde{G}}(\tilde{S})$ then there must be a directed path $v \tuh \dots \tuh w \tuh s$ in $\tilde{G}$ such that $s \in \tilde{S}$, and all other nodes on the directed path are in $V \cup J$.
  The directed path $v \tuh \dots \tuh w$ must also be present in $H$. 
  The node $s$ must be $s_{\{w,u\}}$ for some $u \in V \cup J$ with $w \tut u$ in $H$.
  But since $\dots \suh w \tut u$ cannot occur in a MAG $H$ (because $G$ is acyclic), the directed path must be the trivial path $v = w$.
  Hence $v \tut u$ in $H$, and therefore $v \in \Anc_G(S)$.
  We also state the contra-positive for convenience: for $v \in V \cup J$: if $v \notin \Anc_G(S)$ then $v \notin \Anc_{\tilde{G}}(\tilde{S})$.
  Furthermore, observe that if $x \hus y$ in $H$, then $x \notin \Anc_H(y)$, and therefore $x \notin \Anc_{\tilde{G}}(y)$, and hence $x \notin \Anc_{\tilde{G}}(\{y\} \cup \tilde{S})$.

  It remains to show that $H$ represents $\tilde{G}$ given $\tilde{S}$.
  For that we need to show:
  \begin{enumerate}
    \item For all $x \in V \cup J, y \in V$: $x \sus y$ is in $H$ if and only if there is an inducing path between $x$ and $y$ given $\tilde{S}$ in $\tilde{G}$;
    \item If $x \hus y$ in $H$ then $x \notin \Anc_{\tilde{G}}(\{y\} \cup \tilde{S})$;
    \item If $x \tus y$ in $H$ then $x \in \Anc_{\tilde{G}}(\{y\} \cup \tilde{S})$;
  \end{enumerate}

  First adjacency.
  Let $x \in V \cup J$, $y \in V$.
  Suppose $x$ and $y$ are not adjacent in $H$.
  We must show that in $\tilde{G}$ there is no inducing path between $x$ and $y$ given $\tilde{S}$.
  Suppose there is one.
  Pick one of minimal length.
  It must be a collider path in which every collider is in $\Anc_{\tilde{G}}(\{x,y\} \cup \tilde{S})$.

  Suppose the path contains a node of $\tilde{S}$.
  Then the path has to be of the form $x \dots a \tuh s_{\{a,b\}} \hut b \dots y$ with $a \tut b$ in $H$.
  Since only the endpoint nodes can be non-colliders, the path must be $x \tuh s_{\{x,y\}} \hut y$ with $x \tut y$ in $H$. 
  Contradiction.
  We conclude that the path can only contain nodes in $V \cup J$.

  Assume that each edge on the path is present in $H$ as well (in other words, $i \huh j$ is not on the path).
  Then each collider on the path is not in $\Anc_G(S)$.
  Therefore, by the result above, each collider on the path is not in $\Anc_{\tilde{G}}(\tilde{S})$.
  But this means that each collider on the path must be in $\Anc_{\tilde{G}}(\{x,y\})$, and hence in $\Anc_H(\{x,y\})$.
  So then the path is also an inducing path between $x$ and $y$ in $H$, which contradicts the maximality of $H$.

  Now suppose $i \huh j$ is on the path. 
  All other edges on the path must also be present in $H$.
  The edge $i \huh j$ is not in $H$, but the edge $i \tuh j$ or $i \tut j$ is in $H$.
  %This corresponds with the directed edge $i \tuh j$ in $H$.
  In the latter case, there must be a path of the form $x \suh \dots \huh \dots \huh i \tut j \huh \dots \huh \dots \hus y$ in $H$.
  But this can only happen for $x=i$ and $y=j$, that is, if the path consists of the single edge $i \tut j$.
  But then $x$ would be adjacent to $y$ in $H$, a contradiction.

  Hence the only possibility is that $i \tuh j$ in $H$.
  Without loss of generality (by swapping $x,y$ if necessary), we can then assume the path to be of the form $x \dots i \huh j \dots y$.
  If $x = i$ then we can use the same reasoning as before to show that the path $x = i \tuh j \dots y$ in $H$ must be inducing, contradicting the maximality of $H$.
  If $x \ne i$, then the first part of the inducing path in $\tilde{G}$, say $z_0 \dots z_k$ (with $z_0 = x$ and $z_k = i$ and $k \ge 1$) is also a collider path into $i$ in $H$.
  We will prove by induction that for $1 \le s \le k-1$, $z_s$ is a parent of $j$ in $H$.

  Now suppose that for some $0 \le s \le k-1$, $z_t \tuh j$ is in $H$ for all $t = s+1, \dots, k$ (note this holds in particular for $s=k-1$). 
  We will show that $z_s \tuh j$ is in $H$.
  $z_s$ must be adjacent to $j$ in $H$, otherwise $i \tuh j$ would be visible in $H$.
  The edge between $z_s$ and $j$ cannot be $z_s \hut j$, otherwise $j \tuh z_s \suh z_{s+1} \tuh j$ in $H$, and $H$ would be non-ancestral.  
  It can also not be $z_s \huh j$, otherwise there would be an inducing path $x = z_0 \dots z_s \huh j \dots y$ in $\tilde{G}$ that contains fewer nodes, a contradiction.
  It can also not be $z_s \tut j$, as that would give $i \tuh j \tut z_s$ in $H$ which is not possible.
  Hence it has to be $z_s \tuh j$.

  So by induction, for $1 \le s \le k-1$, $z_s$ is a parent of $j$ in $H$.
  So this gives a collider path $x \suh z_1 \huh \dots \huh z_{k-1} \huh i$ in $H$ such that each $z_s$ is a parent of $j$ in $H$, and $i \tuh j$ in $H$.
  Therefore, $x$ must be adjacent to $j$ in $H$, otherwise $i \tuh j$ would be visible.
  If this edge were of the form $x \hut j$ then $H$ would not be ancestral (because $z_1 \tuh j \tuh x \suh z_1$ in $H$).
  The edge cannot be of the form $x \tut j$ either because then $z_1 \tuh j \tut x$ would be in $H$.
  Hence it must be of the form $x \suh j$.
  But then this edge must also be present in $\tilde{G}$.
  Then we have an inducing path $x \suh j \dots y$ in $\tilde{G}$ that is shorter than the one that was assumed to be of minimal length.
  Contradiction.

  We have shown that if $x$ and $y$ are not adjacent in $H$, there is no inducing path between $x$ and $y$ in $\tilde{G}$ given $\tilde{S}$.
  Now suppose $x$ and $y$ are adjacent in $H$.
  We distinguish four cases corresponding to the type of the edge between $x$ and $y$ in $H$. 
  \begin{itemize}
    \item $x \tuh y$: By construction $x \tuh y$ is in $\tilde{G}$. 
      Then $y \notin \Anc_G(S)$, and therefore also $y \notin \Anc_{\tilde{G}}(\tilde{S})$.
      Therefore, $x \in \Anc_{\tilde{G}}(\{y\} \cup \tilde{S})$ and $y \notin \Anc_{\tilde{G}}(\{y\} \cup \tilde{S})$.
    \item $x \hut y$: This case is similar (with the roles of $x,y$ swapped).
    \item $x \huh y$: By construction $x \huh y$ is in $\tilde{G}$.
      Then both $x, y \notin \Anc_G(S)$, and therefore also $x,y \notin \Anc_{\tilde{G}}(\tilde{S})$. 
      Therefore, $x \notin \Anc_{\tilde{G}}(\{y\} \cup \tilde{S})$ and $y \notin \Anc_{\tilde{G}}(\{y\} \cup \tilde{S})$.
    \item $x \tut y$: Then $x \tuh s_{\{x,y\}} \hut y$ in $\tilde{G}$.
      Therefore, $x \in \Anc_{\tilde{G}}(\{y\} \cup \tilde{S})$ and $y \in \Anc_{\tilde{G}}(\{y\} \cup \tilde{S})$.
  \end{itemize}
  In all cases, there is an inducing path between $x$ and $y$ given $\tilde{S}$ in $\tilde{G}$.
  We conclude that $H$ represents $\tilde{G}$ given $\tilde{S}$.
\end{proof}
\begin{Rem}
  If in addition to the assumptions of Lemma~\ref{lem:zhang2008_lemma_10_mag_extended},
  $H$ contains a directed path from $i$ to $j$ that avoids the edge $i \tuh j$, then there exists such $\tilde{G}$ with the additional property that $i \tuh j$ is not in $\tilde{G}$.
  The proof is almost identical, except that:
  \begin{enumerate}
    \item the directed edge $i \tuh j$ is omitted from $\tilde{G}$;
    \item although not every directed edge in $H$ is present in $\tilde{G}$, the ancestral relations between variables in $V \cup J$ are the same in $H$ as in $\tilde{G}$ by construction;
    \item if $x \tuh y$ in $H$ then $x \tuh y$ is in $\tilde{G}$, except for $i \tuh j$ (for which $i \huh j$ is in $\tilde{G}$ instead).
  \end{enumerate}
\end{Rem}
The results in this subsubsection mean that for directed edges in a MAG, ``invisible'' is synonym for ``possibly confounded'', and ``visible'' is synonym for ``unconfounded''.
This shows that our extension Definition~\ref{def:mag_visible_edge} is appropriate in that it yields the desired properties.
Furthermore, all undirected edges in a MAG are ``possibly confounded''.

Let $G$ be a CADMG with input nodes $J$ and output nodes $V^+ = V \dcup S$. 
Denote the MAG that represents $G$ given $S$ by $H := \MAG(G \given S)$.
We now add an intervention node $I_a$ to $G$, for $a \in V$.
What can we say about the MAG $\MAG(G_{\doit(I_a)} \given S)$?

\begin{Prp}\label{prp:camdg_do_I_mag}
Let $G$ be a CADMG with input nodes $J$ and output nodes $V^+ = V \dcup S$. 
Denote the MAG that represents $G$ given $S$ by $H := \MAG(G \given S)$.
For $A \subseteq V$, define $G^* := G_{\doit(I_A)}$ and $H^* := \MAG(G_{\doit(I_A)} \given S)$.
Then:
  \begin{enumerate}[label=\roman*)]
    \item The subgraph $(H^*)_{V \cup J} = H$.
    \item For all $a \in A$: The edge $I_a \tus a$ is in $H^*$ (note: the edge is of the form $I_a \tuh a$ if $a \notin \Anc_G(S)$, and of the form $I_a \tut a$ if $a \in \Anc_G(S)$).
    \item For all $a \in A, b \in V \sm \{a\}$: if $I_a \sus b \in H^*$, then $I_a \tus b \in H^*$ and $a \tus b \in H$. This implies one of the following mutually exclusive cases:
      \begin{itemize}
        \item $a \tuh b$ is an invisible edge in $H$ and $I_a \tuh b \in H^*$, or
        \item $a \tut b$ in $H$ and $I_a \tut b \in H^*$.
      \end{itemize}
    \item For all $a \in A, b \in V \sm \{a\}$ such that $a \tuh b$ is an invisible edge in $H$, or $a \tut b$ is in $H$: there exists CADMG $\tilde{G}$ with input nodes $J$ and output nodes $\tilde{V}^+ = V \dcup \tilde{S}$ that is represented by $H$ given $\tilde{S}$, and such that $I_a \tus b \in \MAG(\tilde{G}_{\doit(I_a)} \given \tilde{S})$
      (note that $a \tuh b$ in $H$ gives $I_a \tuh b \in \MAG(\tilde{G}_{\doit(I_a)} \given \tilde{S})$, and $a \tut b$ in $H$ gives $I_a \tut b \in \MAG(\tilde{G}_{\doit(I_a)} \given \tilde{S})$).
  \end{enumerate}
\end{Prp}
\begin{proof}
  Note that for $i \in V \cup J \cup S, B \subseteq V \cup J \cup S$, $i \in \Anc_G(B) \iff i \in \Anc_{G*}(B)$.
  \begin{enumerate}[label=\roman*)]
    \item For $i,j \in V \cup J$, $i \in \Anc_G(\{j\} \cup S)$ if and only if $i \in \Anc_{G^*}(\{j\} \cup S)$.
      Furthermore, all non-endpoint nodes on an inducing path must be output nodes.
      Hence, for a pair of nodes in $V \cup J$, an inducing path given $S$ between these nodes in $G$ is an inducing path given $S$ between these nodes in $G^*$, and vice versa.

    \item Let $a \in A$. Note that $I_a \tuh a$ in $G^*$. Hence $I_a \tus a$ in $H^*$. 
      If $a \in \Anc_G(S)$, then $I_a \tut a$ in $H^*$, whereas if $a \notin \Anc_G(S)$, then $I_a \tuh a$ in $H^*$ because $a \notin \Anc_{G^*}(I_a)$. 

    \item Let $a \in A$, $b \in V \sm \{a\}$. 
      If $I_a \sus b \in H^*$ then there exists an inducing path $I_a \tuh a \huh z_1 \huh \dots \huh z_k \huh b$ given $S$ in $G^*$.
      All colliders on this path are in $V \cap \Anc_{G^*}(\{b\} \cup S)$ (since $\Anc_{G^*}(I_a) = \{I_a\}$).
      By omitting the edge $I_a \tuh a$ we obtain an inducing path $a \huh z_1 \huh z_k \huh b$ given $S$ in $G$ (where each collider is in $V \cap \Anc_{G}(\{b\} \cup S)$) that is into $a$.
      Hence, $a$ and $b$ are adjacent in $H$.
  
      \begin{enumerate}
        \item Suppose $a \notin \Anc_G(S)$.
      First, $a \in \Anc_{G^*}(b)$ (because $a \notin \Anc_{G^*}(S)$ since $a \notin \Anc_G(S)$ by assumption), and consequently, $I_a \in \Anc_{G^*}(b)$.
      This also implies $b \notin \Anc_{G^*}(S)$ (because otherwise $a \in \Anc_{G^*}(S)$).
      Hence the edge between $I_a$ and $b$ in $H^*$ must be $I_a \tuh b$.
      In $G$, this implies $a \in \Anc_G(b)$ and $b \notin \Anc_G(S)$.
      Since $G^*$ is acyclic, $b \notin \Anc_{G^*}(a)$ and therefore also $b \notin \Anc_G(a)$.
      Hence $a \tuh b$ in $H$.
      
        \item Suppose $a \in \Anc_G(S)$.
          Then also $a \in \Anc_{G^*}(S)$, and $I_a \in \Anc_{G^*}(S)$.
      Hence the edge between $I_a$ and $b$ in $H^*$ must be $I_a \tus b$.

      Since $a \in \Anc_{G}(S)$, the edge between $a$ and $b$ in $H$ must be $a \tus b$.
      If it is of the form $a \tuh b$, then $b \notin \Anc_{G}(S)$, hence $b \notin \Anc_{G^*}(\{I_a\} \cup S)$, and therefore $I_a \tuh b$ in $H^*$.
      Otherwise (if it is of the form $a \tut b$), then $b \in \Anc_{G}(\{a\} \cup S)$, hence $b \in \Anc_{G^*}(\{a\} \cup S)$, and since $a \in \Anc_{G^*}(S)$, surely $b \in \Anc_{G^*}(S)$.
      Hence, also $I_a \tut b$ in $H^*$.
      \end{enumerate}
      If $a \tuh b$ is in $H$, then it must be invisible, because the inducing path in $G$ given $S$ between $a$ and $b$ is into $a$ (Lemma~\ref{lem:zhang2008_lemma_9_mag_extended}).
    \item Let $a \in A$, $b \in V \sm \{a\}$. 
      Suppose that $a \tuh b$ is an invisible edge in $H$, or $a \tut b$ in $H$.
      Then we can consider $\tilde{G}$ as constructed in the proof of Lemma~\ref{lem:zhang2008_lemma_10_mag_extended}.
      This contains the bidirected edge $a \huh b$ and is represented by $H$.
      The path $I_a \tuh a \huh b$ in $\tilde{G}^*$ is inducing given $\tilde{S}$.
      If $a \tuh b$ is in $H$, then $I_a \tuh b$ must be in $\MAG(\tilde{G}_{\doit(I_a)}|\tilde{S})$,
      whereas if $a \tut b$ is in $H$, then $I_a \tut b$ must be in $\MAG(\tilde{G}_{\doit(I_a)}|\tilde{S})$.
  \end{enumerate}
\end{proof}
So we can define:
\begin{Def}\label{def:mag_do_I}
  Let $H$ be a MAG with input nodes $J$ and output nodes $V$, that represents some CADMG with input nodes $J$ and output nodes $V^+ = V \dcup S$ given $S$.
  Let $A \subseteq V$.
  Define $H_{\doit(I_A)}$ as the PAG with input nodes $J \dcup \{I_a : a \in A\}$ and output nodes $V$, such that the subgraph $(H_{\doit(I_A)})_{V \cup J}$ equals $H_{V \cup J}$, and in addition has edges
  \begin{itemize}
%    \item $I_a \tuo a$ for all $a \in A$ 
    \item $I_a \tuh a$ if $a$ has an adjacent arrowhead in $H$, $I_a \tut a$ if $a$ has an adjacent undirected edge in $H$, and $I_a \tuo a$ otherwise;
    \item $I_a \tuh b$ for all $a \in A, b \in V \sm \{a\}$ such that $a \tuh b$ is an invisible edge in $H$;
    \item $I_a \tut b$ for all $a \in A, b \in V \sm \{a\}$ such that $a \tut b$ is in $H$.
  \end{itemize} 
\end{Def}
Note that we loose information about confounding by going to the MAG level: in general, $\MAG(G_{\doit(I_a)}|S) \ne \MAG(G|S)_{\doit(I_a)}$.
But we do have that $\MAG(G_{\doit(I_a)}|S) \subseteq (\MAG(G|S))_{\doit(I_a)}$.
This also means we may have lost conditional independence information.
Still, the operation preserves as much conditional independence information as possible.
We can use it to obtain a causal interpretation for MAGs (in the presence of selection bias), derive a do-calculus for MAGs, derive adjustment criteria, and an ID algorithm.
\Joris{For example, would the natural definition of ``$a$ causes $b$ in the \emph{selected} population'' be: $b \nIndep I_a \given a$?}
\Joris{A generalization to PAGs might work as follows.
The first point can remain the same.
The condition for the second point should probably be weakened to: if $a \tuh b$ is a possibly invisible edge in $H$ (that is, if there exists a MAG represented by the PAG in which $a \tuh b$ is an invisible edge).
Hm, if $a \sus b$ in $H$, then it can be:
$a \tut b$ (and then we add $I_a \tut b$), $a \tuh b$ (and then we add $I_a \tuh b$ if $a \tuh b$ is invisible), $a \hut b$ (then we don't add $I_a \sus b$), $a \huh b$ (then we don't add $I_a \sus b$).
So for $a \ouo b$ in $H$ we could add $I_a \tuo b$,
for $a \ouh b$ we add $I_a \tuh b$, 
for $a \huo b$ we don't add an edge between $I_a$ and $b$,
for $a \out b$ we add $I_a \tut b$ if $b$ has no adjacent arrowheads,
for $a \tuo b$ we add $I_a \tuo b$.
This may not yet be optimally informative, though (sound but not necessarily complete).
The third point should probably be weakened to: if $a \tut b$ might be in $H$ (that is, if there exists a MAG represented by the PAG that contains $a \tut b$).
It's not clear whether, if we start with a valid MAG, the operation of adding an intervention node commutes with itself (for disjoint targets).
Note that ``definitely visible'' in a PAG is the same as ``visible'' in a PAG.
``Possibly invisible'' means that there exists a MAG in which it is ``invisible'', that is in which it is not ``visible''.
Is that equivalent to ``not definitely visible''?
}
One might hope that given a MAG $H$ that represents a CADMG $G$, the PAG $H_{\doit(I_A)}$ represents $G_{\doit(I_A)}$ given $S$.
This does not hold in general, as the example in Figure~\ref{fig:mag_do_I} shows.
\begin{figure}\centering
  \begin{tikzpicture}
    \begin{scope}
      \node at (0.75,0.75) {$G^1$:};
      \node[ndout] (X2) at (0,0) {$a$};
      \node[ndout] (X3) at (1.5,0) {$b$};
      \draw[tuh] (X2) edge (X3);
    \end{scope}
    \begin{scope}[xshift=5cm]
      \node at (0,0.75) {$G^1_{\doit(I_a)}$:};
      \node[ndout] (X1) at (-1.5,0) {$I_a$};
      \node[ndout] (X2) at (0,0) {$a$};
      \node[ndout] (X3) at (1.5,0) {$b$};
      \draw[tuh] (X1) edge (X2);
      \draw[tuh] (X2) edge (X3);
    \end{scope}
    \begin{scope}[xshift=10cm]
      \node at (0,0.75) {$\MAG(G^1_{\doit(I_a)})$:};
      \node[ndout] (X1) at (-1.5,0) {$I_a$};
      \node[ndout] (X2) at (0,0) {$a$};
      \node[ndout] (X3) at (1.5,0) {$b$};
      \draw[tuh] (X1) edge (X2);
      \draw[tuh] (X2) edge (X3);
    \end{scope}
    \begin{scope}[yshift=-2cm]
      \node at (0.75,0.75) {$G^2$:};
      \node[ndout] (X2) at (0,0) {$a$};
      \node[ndout] (X3) at (1.5,0) {$b$};
      \draw[tuh] (X2) edge (X3);
      \draw[huh,bend left] (X2) edge (X3);
    \end{scope}
    \begin{scope}[xshift=5cm,yshift=-2cm]
      \node at (0,0.75) {$G^2_{\doit(I_a)}$:};
      \node[ndout] (X1) at (-1.5,0) {$I_a$};
      \node[ndout] (X2) at (0,0) {$a$};
      \node[ndout] (X3) at (1.5,0) {$b$};
      \draw[tuh] (X1) edge (X2);
      \draw[tuh] (X2) edge (X3);
      \draw[huh,bend left] (X2) edge (X3);
    \end{scope}
    \begin{scope}[xshift=10cm,yshift=-2cm]
      \node at (0,0.75) {$\MAG(G^2_{\doit(I_a)})$:};
      \node[ndout] (X1) at (-1.5,0) {$I_a$};
      \node[ndout] (X2) at (0,0) {$a$};
      \node[ndout] (X3) at (1.5,0) {$b$};
      \draw[tuh] (X1) edge (X2);
      \draw[tuh] (X2) edge (X3);
      \draw[tuh,bend right] (X1) edge (X3);
    \end{scope}
    \begin{scope}[yshift=-4.5cm]
      \node at (0.75,0.75) {MAG $H$:};
      \node[ndout] (X2) at (0,0) {$a$};
      \node[ndout] (X3) at (1.5,0) {$b$};
      \draw[tuh] (X2) edge (X3);
    \end{scope}
    \begin{scope}[xshift=10cm,yshift=-4.5cm]
      \node at (0,0.75) {PAG $H_{\doit(I_a)}$:};
      \node[ndout] (X1) at (-1.5,0) {$I_a$};
      \node[ndout] (X2) at (0,0) {$a$};
      \node[ndout] (X3) at (1.5,0) {$b$};
      \draw[tuo] (X1) edge (X2);
      \draw[tuh] (X2) edge (X3);
      \draw[tuh,bend right] (X1) edge (X3);
    \end{scope}
    \begin{scope}[yshift=-7cm]
      \node at (0.75,0.75) {$G^3$:};
      \node[ndout] (X2) at (0,0) {$a$};
      \node[ndout] (X3) at (1.5,0) {$b$};
      \node[ndout] (XS) at (0,-1) {$s$};
      \draw[tuh] (X2) edge (X3);
      \draw[tuh] (X2) edge (XS);
    \end{scope}
    \begin{scope}[xshift=5cm,yshift=-7cm]
      \node at (0,0.75) {$G^3_{\doit(I_a)}$:};
      \node[ndout] (X1) at (-1.5,0) {$I_a$};
      \node[ndout] (X2) at (0,0) {$a$};
      \node[ndout] (X3) at (1.5,0) {$b$};
      \node[ndout] (XS) at (0,-1) {$s$};
      \draw[tuh] (X1) edge (X2);
      \draw[tuh] (X2) edge (X3);
      \draw[tuh] (X2) edge (XS);
    \end{scope}
    \begin{scope}[xshift=10cm,yshift=-7cm]
      \node at (0,0.75) {$\MAG(G^3_{\doit(I_a)})$:};
      \node[ndout] (X1) at (-1.5,0) {$I_a$};
      \node[ndout] (X2) at (0,0) {$a$};
      \node[ndout] (X3) at (1.5,0) {$b$};
      \draw[tut] (X1) edge (X2);
      \draw[tuh] (X2) edge (X3);
    \end{scope}
    \begin{scope}[yshift=-10cm]
      \node at (0.75,0.75) {$G^4$:};
      \node[ndout] (X2) at (0,0) {$a$};
      \node[ndout] (X3) at (1.5,0) {$b$};
      \node[ndout] (XS) at (0,-1) {$s$};
      \draw[tuh] (X2) edge (X3);
      \draw[huh,bend left] (X2) edge (X3);
      \draw[tuh] (X2) edge (XS);
    \end{scope}
    \begin{scope}[xshift=5cm,yshift=-10cm]
      \node at (0,0.75) {$G^4_{\doit(I_a)}$:};
      \node[ndout] (X1) at (-1.5,0) {$I_a$};
      \node[ndout] (X2) at (0,0) {$a$};
      \node[ndout] (X3) at (1.5,0) {$b$};
      \node[ndout] (XS) at (0,-1) {$s$};
      \draw[tuh] (X1) edge (X2);
      \draw[tuh] (X2) edge (X3);
      \draw[huh,bend left] (X2) edge (X3);
      \draw[tuh] (X2) edge (XS);
    \end{scope}
    \begin{scope}[xshift=10cm,yshift=-10cm]
      \node at (0,0.75) {$\MAG(G^4_{\doit(I_a)})$:};
      \node[ndout] (X1) at (-1.5,0) {$I_a$};
      \node[ndout] (X2) at (0,0) {$a$};
      \node[ndout] (X3) at (1.5,0) {$b$};
      \draw[tut] (X1) edge (X2);
      \draw[tuh] (X2) edge (X3);
      \draw[tuh,bend right] (X1) edge (X3);
    \end{scope}
  \end{tikzpicture}
  \caption{MAG $H$ represents two different CADMGs $G^1, G^2$.
    After adding an intervention node $I_a$ to the CADMGs, their representative MAGs no longer coincide.
    PAG $H_{\doit(I_a)}$ represents $\MAG(G^2_{\doit(I_a)})$, but not $\MAG(G^1_{\doit(I_a)})$.
    MAG $H$ also represents two different CADMGs $G^3, G^4$ given $\{s\}$.
    After adding an intervention node $I_a$ to the CADMGs, their representative MAGs no longer coincide.
    PAG $H_{\doit(I_a)}$ represents $\MAG(G^4_{\doit(I_a)})$, but not $\MAG(G^3_{\doit(I_a)})$.
    \label{fig:mag_do_I}}
\end{figure}
Whereas $H_{\doit(I_A)}$ has too many edges in general to represent $G_{\doit(I_A)}$ (given $S$), it does still encode the ancestral relations and separations of $G_{\doit(I_A)}$ (given $S$).
\Joris{We can formalize this type of representation as follows.
\begin{Def}
  Given a PAG-like graph $H$ (not necessarily valid), and a valid MAG $M$, we say that $H$ \emph{weakly represents a MAG $M$} if $H$ has the same nodes as $M$, the skeleton of $M$ is contained in the skeleton of $H$, and for every $i,j$ adjacent in $M$, each arrowhead/tail edge mark present on the edge $i \sus j$ in $H$ must also be present on the corresponding edge $i \sus j$ in $M$.
  Given an ADMG $G$ and a subset $S$ of selection variables, we say that $H$ weakly represents $G$ given $S$ if $H$ weakly represents $MAG(G \given S)$.
\end{Def}
This should be equivalent to: there exists an orientation $K$ of $H$ (way of replacing each circle edge mark by an arrowhead or a tail) such that $M$ is a subgraph of $K$.
I believe that we then can conclude that $H_{\doit(I_A)}$ weakly represents all $G_{\doit(I_A)}$ given $S$ ( for $G$ represented by $H$ given $S$), and $H_{\doit(D)}$ weakly represents all possible $G_{\doit(D)}$ given $S$ (for $G$ represented by $H$ given $S$).
Alternatively: we could also refer to the ``PAG-like graph'' $H_{\doit(I_A)}$ as a ``SPAG'', a ``superPAG''!}

Note that the PAG $H_{\doit(I_A)}$ is as informative about the result of $\doit(I_A)$ on CADMGs as it could be.
\begin{Prp}
  Let $H$ be a MAG with input nodes $J$ and output nodes $V$, that represents some CADMG with input nodes $J$ and output nodes $V^+ = V \dcup S$ given $S$.
  Let $A \subseteq V$.
  Then for each possible orientation $H^*$ of $H_{\doit(I_A)}$ (obtained by orienting each circle in $H_{\doit(I_A)}$ into a tail or arrowhead), there exists a 
  % PAG $H_{\doit(I_a)}$ represents $G_{\doit(I_a)}$ given $S$   THIS DOES NOT HOLD!
  a CADMG $\tilde{G}$ with input $J$ and output nodes $V^+ = V \dcup \tilde{S}$ with $\MAG(\tilde{G} \given \tilde{S}) = H$ and  $\MAG(\tilde{G}_{\doit(I_A)} \given \tilde{S}) = H^*$.
\end{Prp}
\begin{proof}
  The circles in $H_{\doit(I_A)}$ are on edges of the type $I_a \tuo a$ for nodes $a$ which have no adjacent arrowhead nor an adjacent undirected edge.
  Let $H^*$ be a possible orientation $H^*$ of $H_{\doit(I_A)}$. 
  This amounts to deciding for each $a \in A$ whether $a$ will become ancestor of a selection variable, or not.
  We can then pick $\tilde{G}$ as constructed in Lemma~\ref{lem:representing_CADMG_of_MAG} and see that it does the job.
\end{proof}

In the following, write $I_D := \{I_d : d \in B\}$.
\begin{Prp}
  Let $H$ be a MAG with input nodes $J$ and output nodes $V$, that represents some CADMG with input nodes $J$ and output nodes $V^+ = V \dcup S$ given $S$.
  Let $D \subseteq V$.
  Let $A, B, C \subseteq V \cup J$.
  Let $K$ be a MAG that is contained in $H_{\doit(I_D)}$.
  Then there is a definite $m$-connecting path from $A$ to $B$ given $C$ in $H_{\doit(I_D)}$ if and only if there is a $m$-connecting path from $A$ to $B$ given $C$ in $K$.
\end{Prp}
\begin{proof}
  Suppose there is a $m$-connecting path from $A$ to $B$ given $C$ in some MAG contained in $H_{\doit(I_D)}$.
  The only possible edges on the path for which the corresponding edge in $H_{\doit(I_D)}$ contains a circle are the edges $I_d \tuh d$ with $d$ having no adjacent arrowhead in $H$; in that case, $d$ can only be a non-collider on the path.
  On the corresponding path in $H_{\doit(I_D)}$ these nodes can only appear as definite non-colliders.
  All other nodes on the corresponding path in $H_{\doit(I_D)}$ have no adjacent circle edge marks and hence they are a definite (non-)collider if and only if they are a (non-)collider in the MAG.
  Hence the corresponding path in $H_{\doit(I_D)}$ must be definitely $m$-connecting from $A$ to $B$ given $C$.

  Vice versa, suppose there is a definite $m$-connecting path from $A$ to $B$ given $C$ in $H_{\doit(I_D)}$.
  Then in every MAG contained in $H_{\doit(I_D)}$ (and hence in $K$), the corresponding path is $m$-connecting from $A$ to $B$ given $C$.
\end{proof}
This implies:
\begin{Cor}
  Let $H$ be a MAG with input nodes $J$ and output nodes $V$, that represents some CADMG with input nodes $J$ and output nodes $V^+ = V \dcup S$ given $S$.
  Let $D \subseteq V$.
  Let $A, B, C \subseteq V \cup J$.
  \begin{enumerate}[label=\alph*)]
    \item If there is a definite $m$-connecting path from $A$ to $B$ given $C$ in $H_{\doit(I_D)}$, there exists a CADMG $\tilde{G}$ with input $J$ and output nodes $V^+ = V \dcup \tilde{S}$ such that $\tilde{G}$ is represented by $H$ given $\tilde{S}$, and in $\tilde{G}_{\doit(I_D)}$ there exists a $d$-open path from $A$ to $B$ given $C \cup \tilde{S}$.
    \item If there is no definite $m$-connecting path from $A$ to $B$ given $C$ in $H_{\doit(I_D)}$, then for all CADMGs $\tilde{G}$ with input $J$ and output nodes $V^+ = V \dcup \tilde{S}$ such that $\tilde{G}$ is represented by $H$ given $\tilde{S}$, there is no $d$-connecting path from $A$ to $B$ given $C \cup \tilde{S}$ in $\tilde{G}_{\doit(I_D)}$.
  \end{enumerate}
\end{Cor}
\begin{proof}
  \begin{enumerate}
    \item Let $K$ be a MAG contained in $H_{\doit(I_D)}$ \Joris{Does it exist? Yes, any possible orientation gives a valid MAG (see earlier proposition)}. 
      There exists a CADMG $\tilde{G}$ with input $J$ and output nodes $V^+ = V \dcup \tilde{S}$ with $\MAG(\tilde{G} \given \tilde{S}) = H$ and $\MAG(\tilde{G}_{\doit(I_A)} \given \tilde{S}) = K$.
      If there is a definite $m$-connecting path from $A$ to $B$ given $C$ in $H_{\doit(I_D)}$, there is an $m$-connecting path from $A$ to $B$ given $C$ in $K$,
      and hence there is a $d$-open path from $A$ to $B$ given $C \cup \tilde{S}$ in $\tilde{G}_{\doit(I_A)}$.
\Joris{One could even show: If there is a definite $m$-connecting path from $A$ to $B$ given $C$ in $H_{\doit(I_D)}$, then for any MAG $K$ contained in $H_{\doit(I_D)}$, there exists
  a CADMG $\tilde{G}$ and set $\tilde{S}$ such that $\tilde{G}$ is represented by $H$ given $\tilde{S}$, and $\tilde{G}_{\doit(I_D)}$ is represented by $K$ given $\tilde{S}$, and there exists a $d$-open path in $\tilde{G}_{\doit(I_D)}$ from $A$ to $B$ given $C$.}
    \item We prove the converse.
      Suppose there is a CADMG $\tilde{G}$ with input $J$ and output nodes $V^+ = V \dcup \tilde{S}$ such that $\tilde{G}$ is represented by $H$ given $\tilde{S}$, and with a $d$-connecting path from $A$ to $B$ given $C \cup \tilde{S}$ in $\tilde{G}_{\doit(I_D)}$.
      Then $H^* := \MAG(\tilde{G}_{\doit(I_D)} \given \tilde{S})$ is a subgraph of an orientation $K$ of $H_{\doit(I_D)}$ \Joris{TODO???}.
      Since there is a $d$-connecting path from $A$ to $B$ given $C \cup \tilde{S}$ in $\tilde{G}_{\doit(I_D)}$,
      there is an $m$-connecting path from $A$ to $B$ given $C$ in $H^*$.
      The same path must be present in $K$ and be $m$-connecting given $C$ in $K$ \Joris{TODO???}.
      Hence, there is a definite $m$-connecting path from $A$ to $B$ given $C$ in $H_{\doit(I_D)}$.
  \end{enumerate}
\end{proof}

\Joris{Leihao shows that the notion of (in)visible edge in $H$ is preserved in $H_{\doit(I_a)}$.
Can we show this by using that it is synonym with ``possibly confounded'' in the represented CADMGs?}

\subsubsection{Modeling hard interventions with MAGs}

\Joris{This subsubsection is new and still completely speculative.}

Let $G$ be a CADMG with input nodes $J$ and output nodes $V^+ = V \dcup S$. 
Denote the MAG that represents $G$ given $S$ by $H := \MAG(G \given S)$.
We now perform hard interventions $\doit(D)$ on $G$, for intervention targets $D \subseteq V \cup J$.
What can we infer about $\MAG(G_{\doit(D)} \given S)$ given only $H$?

We define:
\begin{Def}\label{def:do_mag}
  Let $H$ be a MAG with input nodes $J$ and output nodes $V$, that represents some CADMG with input nodes $J$ and output nodes $V^+ = V \dcup S$ given $S$.
  Let $D \subseteq V \cup J$.
  Then $H_{\bar{D}}$ is the mixed graph with input nodes $J \cup D$, output nodes $V \sm D$ and as the edges in $H$ except those that are into $D$.
\end{Def}
\begin{Def}
  For a mixed graph $H$, we denote by $\Un_H$ the set of vertices of $H$ with no adjacent arrowheads.
\end{Def}
We can show some elementary properties:
\begin{Lem}
  Let $G$ be a CADMG with input nodes $J$ and output nodes $V^+ = V \dcup S$. 
  Denote the MAG that represents $G$ given $S$ by $H := \MAG(G \given S)$.
  Let $D \subseteq V \cup J$.
  Define $G^* := G_{\doit(D)}$ and $H^* := \MAG(G^* \given S)$.
  Then:
  \begin{enumerate}
    \item if vertices $i,j$ are adjacent in $H^*$ then they are adjacent in $H_{\bar D}$;
    \item if $i \tus j$ in $H^*$ then $i \tus j$ in $H_{\bar D}$;
  \end{enumerate}
  \Joris{Additionally, Leihao shows that under the assumption that each $d \in D$ has an adjacent arrowhead in $H$:
  if $i \suh j$ in $H^*$ then $i \tut j$ not in $H_{\bar D}$. 
  This means that the differences between $H^*$ and $H_{\bar D}$ are that $H^*$ may have fewer edges than $H_{\bar D}$, and some bidirected edges in $H^*$ may be directed in $H_{\bar D}$.
  
  He further shows that: if $i \tuh j$ in $H_{\bar D}$ is visible (in $H_{\bar D}$) then $i \tuh j$ must be visible in $H$;
  and hence (see old stuff below) $i \tuh j$ must be in $H^*$.

  I note myself that the Lemma in footnote 22 can be applied to our setting (an invisible directed edge $a \tuh b$ in $H_{\bar D}$ that is bidirected in $H^*$) because:
  \begin{itemize}
    \item (speculative) if there is another directed path from $a$ to $b$ in $H_{\bar D}$ other than $a \tuh b$ then perhaps we can recurse?
    \item this follows from the assumpion that $a \tuh b$ is invisible
    \item this also follows from the assumption that $a \tuh b$ is invisible: such a discriminating path would either end as $q_n \hut a \tuh b$ with $q_n \tuh d$ (violating t1), or as $q_n \huh a \tuh b$ (violating the invisibility assumption)
  \end{itemize}
  Still we have to understand why making this change won't make another `ghost edge' visible!
  }
\end{Lem}
\begin{proof}
  Note that $x \in \Anc_{G_{\doit(D)}}(U)$ implies that $x \in \Anc_G(U)$, for all $x \in V^+ \cup J$ and $U \subseteq V^+ \cup J$.
  \begin{enumerate}
    \item
      Let $i \sus j$ be an edge in $H^*$. 
      Then there is an inducing path $\pi$ between $i$ and $j$ given $S$ in $G_{\doit(D)}$.
      Every non-endpoint node on the path must be a collider in $\Anc_{G_{\doit(D)}}(\{i,j\} \cup S)$.
      The same sequence of node indices defines a path in $G$ (with possibly different node types), where every non-endpoint node must be a collider in $\Anc_G(\{i,j\} \cup S)$.
      Therefore, there is an inducing path between $i$ and $j$ given $S$ in $G$.
      Hence, $i$ and $j$ are adjacent in $H$.
      
      If $H^*$ had an edge $i \suh d$ for $d \in D$ then in $G^*$ there must be an inducing path between $i$ and $d$ given $S$ that is \emph{into} $d$.
      However, $G^*$ contains no edges into $d$.
      Contradiction.
      So $H^*$ has no edges into $D$.
    \item
      Suppose the edge $i \sus j$ in $H^*$ were out of $i$.
      Then $i \in \Anc_{G_{\doit(D)}}(\{j\} \cup S)$, and hence $i \in \Anc_G(\{j\} \cup S)$.
      Therefore, the edge $i \sus j$ in $H_{\bar d}$ must be out of $i$ as well.
      Similarly, if the edge $i \sus j$ in $H^*$ were out of $j$, the edge $i \sus j$ in $H_{\bar D}$ must be out of $j$ as well.
  \end{enumerate}
\end{proof}
We extend Lemma 12 of \cite{Zhang2008_AI} to input nodes and selection bias.
\begin{Lem}\label{lem:zhang_2008_lemma_12}
  Let $H$ be a MAG with input nodes $J$ and output nodes $V$, that represents some CADMG with input nodes $J$ and output nodes $V^+ = V \dcup S$ given $S$.
  Let $D \subseteq V \cup J$.
  Define $G^* := G_{\doit(D)}$ and $H^* := \MAG(G^* \given S)$.
  For $a, b \in V \cup J$ and $C \subseteq V \cup J$:
  If there is an $m$-connecting path between $a$ and $b$ given $C$ in $H^*$ then there is one in $H_{\bar{D}}$.
\end{Lem}
\begin{proof}
  \cite{Zhang2008_AI} proves the special case $J = \emptyset$ and $S = \emptyset$.
  The generalization to $J \ne \emptyset$ is trivial since the node type does not matter in the definition of $m$-connecting path.
  Let $\pi$ be an $m$-connecting path between $a$ and $b$ given $C$ in $H^*$.
  We can decompose it into segments consisting solely of undirected edges and segments containing no undirected edges.
  The undirected edges on $\pi$ must also be present in $H_{\bar{D}}$ (see previous Lemma).
  All vertices adjacent to an undirected edge on $\pi$ are non-colliders, and hence, cannot be in $C$.
  So the undirected segments are also $m$-connecting paths in $H_{\bar{D}}$.

  The subgraph $K$ of $H^*$ obtained by omitting the undirected edges is a DMAG.
  Indeed, as a subgraph of MAG $H^*$ and without undirected edges we only have to show that it is maximal.
  Suppose there is an inducing path between two nodes in $K$.
  We have to show that the two nodes are adjacent.
  The inducing path in $K$ must also exist in $H^*$.
  Furthermore, it must also be inducing in $H^*$.
  Because $H^*$ is maximal, the two nodes are adjacent in $H^*$.
  We will show that the edge between them cannot be undirected, and must therefore be present in $K$ as well.
  Indeed, if the inducing path in $K$ consists of a single edge, it cannot be undirected.
  If it consists of multiple edges, then it contains a collider that is ancestor of one of the two nodes, and hence that node has an adjacent arrowhead in $K$ (and in $H^*$), which contradicts the possibility of an adjacent undirected edge in $H^*$.

  Note that the same ancestral relations hold in $K$ as in $H^*$.
  Therefore, every segment on $\pi$ that is contained in $K$ is $m$-connecting in $H^*$ and hence in $K$.
  So we can find a corresponding $m$-connecting path between the endpoints of the segment in $H_{\bar{D}}$ according to \cite{Zhang2008_AI}'s Lemma 12.
  \Joris{Really? We have to relate $K$ with some $K_{\bar{D}}$ and show that $K_{\bar{D}}$ is a subgraph of $H_{\bar{D}}$?
  Lemma 12 would say: if there is an $m$-connecting path between $a$ and $b$ given $C$ in a DMAG $\MAG(G_{\doit(D)})$ then there is one in $\MAG(G)_{\doit(D)}$.
  Perhaps we can show that $K$ is (up to node types) equal to $(H^*)_{\doit(D,\Un_{H^*})}$? Or to $\MAG(G_{\doit(D,\Un_{H^*})} \given S)$?
  Wait, this is circular!
  It would be really nice if we can show that $(H^*)_{\overline{\Un_{H^*}}} = \MAG(G_{\doit(D,\Un_{H^*})} \given S)$.

  Perhaps we need to define $E := D \cup \Un_{H^*}$ and consider $G_{\doit(E)}$ and its $\MAG(G_{\doit(E)} \given S)$?
  Can we show that its MAG equals $(H^*)_{\bar{\Un_{H^*}}}$?
  First, can we show that $\Un_{H^*} = \Un_H$?
  No, we can get new arrowheads in $H^*$!

  Aha. Leihao didn't mean to apply Lemma 12, but only the Lemma in the footnote on page 22.
  }
  We can now paste all these segments together to an $m$-connecting path $\tilde{\pi}$ in $H_{\bar{D}}$.
  Note that the nodes where we connect segments are non-colliders on $\pi$ and remain non-colliders on $\tilde{\pi}$, and hence do not block $\tilde{\pi}$.
\end{proof}

\begin{figure}\centering
  \begin{tikzpicture}
    \begin{scope}[xshift=-4cm]
      \node[ndout] (a) at (0,0) {$a$};
      \node[ndout] (b) at (0,-1.5) {$b$};
      \node[ndout] (z) at (-1.5,0) {$z$};
      \node[ndout] (d1) at (-0.75,1) {$d_1$};
      \node[ndout] (d2) at (1,-0.75) {$d_2$};
      \draw[huh] (a) edge (z);
      \draw[tuh] (a) edge (d1);
      \draw[tuh] (d1) edge (z);
      \draw[huh] (a) edge (b);
      \draw[tuh] (a) edge (d2);
      \draw[tuh] (d2) edge (b);
      \draw[tuh] (z) edge (b);
      \node at (0,-2.5) {$G$ {\color{red}[$\C{G}$]}};
    \end{scope}
    \begin{scope}[xshift=0cm]
      \node[ndout] (a) at (0,0) {$a$};
      \node[ndout] (b) at (0,-1.5) {$b$};
      \node[ndout] (z) at (-1.5,0) {$z$};
      \node[ndout] (d1) at (-0.75,1) {$d_1$};
      \node[ndout] (d2) at (1,-0.75) {$d_2$};
      \draw[tuh] (a) edge (z);
      \draw[tuh] (a) edge (d1);
      \draw[tuh] (d1) edge (z);
      \draw[tuh] (a) edge (b);
      \draw[tuh] (a) edge (d2);
      \draw[tuh] (d2) edge (b);
      \draw[tuh] (z) edge (b);
      \draw[tuh,red] (d1) edge (b);
      \node at (0,-2.5) {$H = \MAG(G)$ {\color{red}[$\C{M}$]}};
    \end{scope}
    \begin{scope}[xshift=4cm]
      \node[ndout] (a) at (0,0) {$a$};
      \node[ndout] (b) at (0,-1.5) {$b$};
      \node[ndout] (z) at (-1.5,0) {$z$};
      \node[ndint] (d1) at (-0.75,1) {$d_1$};
      \node[ndint] (d2) at (1,-0.75) {$d_2$};
      \draw[huh] (a) edge (z);
      \draw[tuh] (d1) edge (z);
      \draw[huh] (a) edge (b);
      \draw[tuh] (d2) edge (b);
      \draw[tuh] (z) edge (b);
      \node at (0,-2.5) {$G_{\doit(X)}$ {\color{red}[$\C{G}_{\underline{Y}\overline{X}}$]}};
    \end{scope}
    \begin{scope}[xshift=8cm]
      \node[ndout] (a) at (0,0) {$a$};
      \node[ndout] (b) at (0,-1.5) {$b$};
      \node[ndout] (z) at (-1.5,0) {$z$};
      \node[ndint] (d1) at (-0.75,1) {$d_1$};
      \node[ndint] (d2) at (1,-0.75) {$d_2$};
      \draw[huh] (a) edge (z);
      \draw[tuh] (d1) edge (z);
      \draw[huh] (a) edge (b);
      \draw[tuh] (d2) edge (b);
      \draw[tuh] (z) edge (b);
      \node at (0,-2.5) {$H^* = \MAG(G_{\doit(X)})$ {\color{red}[$\C{M}_{\C{G}_{\underline{Y}\overline{X}}}$]}};
    \end{scope}
    \begin{scope}[yshift=-5cm,xshift=-4cm]
      \node[ndout] (a) at (0,0) {$a$};
      \node[ndout] (b) at (0,-1.5) {$b$};
      \node[ndout] (z) at (-1.5,0) {$z$};
      \node[ndint] (d1) at (-0.75,1) {$d_1$};
      \node[ndint] (d2) at (1,-0.75) {$d_2$};
      \draw[tuh] (a) edge (z);
      \draw[tuh] (d1) edge (z);
      \draw[tuh] (a) edge (b);
      \draw[tuh] (d2) edge (b);
      \draw[tuh] (z) edge (b);
      \draw[tuh,red] (d1) edge (b);
      \node at (0,-2.5) {$H_{\bar X}$ {\color{red}[$\C{M}_{\underline{Y}\overline{X}}$]}};
    \end{scope}
    \begin{scope}[yshift=-5cm,xshift=0cm]
      \node[ndout] (a) at (0,0) {$a$};
      \node[ndout] (b) at (0,-1.5) {$b$};
      \node[ndout] (z) at (-1.5,0) {$z$};
      \node[ndint] (d1) at (-0.75,1) {$d_1$};
      \node[ndint] (d2) at (1,-0.75) {$d_2$};
      \draw[huh] (a) edge (z);
      \draw[tuh] (d1) edge (z);
      \draw[tuh] (a) edge (b);
      \draw[tuh] (d2) edge (b);
      \draw[tuh] (z) edge (b);
      \draw[tuh,red] (d1) edge (b);
      \node at (0,-2.5) {$(H_{\bar X})'$ {\color{red}[$\C{M}_{\underline{Y}\overline{X}}'$]}};
    \end{scope}
  \end{tikzpicture}
  \caption{Counterexample to Zhang's claim: ``Furthermore, making this change will not make any other such directed edge in $\C{M}_{\underline{Y}\overline{X}}$ [$H_{\bar X}$] visible.''
  $X = \{d_1,d_2\}$. $a \tuh z$ and $a \tuh b$ in $H_{\bar X}$ are both invisible (and are $a \huh z$ and $a \huh b$ in $H^*$).
  $(H_{\bar X})'$: Replacing $a \tuh z$ by $a \huh z$ in $H_{\bar X}$ makes $a \tuh b$ visible!
  Also, after the replacement, $d_1 \tuh z \huh a \tuh b$ is a discriminating path for $a$.
  Note that $H_{\bar X}$ and $H^*$ are not even Markov equivalent! ($d_1 \mPerp b \given a, z$ in $H_{\bar X}$, but $d_1 \nmPerp b \given a, z$ in $H^*$)
  \Joris{Ah... Leihao found my mistake: I overlooked the inducing path $d_1 \tuh z \huh a \huh b$ in $G$, and hence the directed edge $d_1 \tuh b$ in $H$.}}
\end{figure}
Now the final step is to combine adding intervention nodes with performing hard interventions.
\Joris{TODO. 
This seems a bit problematic since we first would like to perform a hard intervention, and then add intervention nodes.
But if we perform a hard intervention, we do not end up with a PAG that can be seen as a set of possible MAGs to which we can then add intervention nodes.
So perhaps we want to turn around the ordering?
We could, since on $G$ we perform $G_{\doit(D,I_A)}$ and the ordering is arbitrary since $D$ and $A$ are disjoint.
If they are not disjoint, then first $\doit(d)$ and then $\doit(I_d)$ is equivalent to just $\doit(d)$;
while first $\doit(I_d)$ and then $\doit(d)$ is equivalent to adding an isolated node $I_d$ combined with $\doit(d)$.
On the MAG/PAG level, if we first $\doit(I_d)$ we may get additional edge $I_d \tuh b$ for example if $d \tuh b$
is invisible; if we then $\doit(b)$ then this edge disappears again, which suggests the ordering doesn't matter.
Suppose we first $\doit(I_d)$ and get $I_d \tut b$ if $d \tut b$ is present; if we then $\doit(b)$ the edges
$d \tut b$ and $I_d \tut b$ stay. If we turn it around, it also works. So graphically, it seems the ordering 
makes no difference. Interestingly, if we first $\doit(I_D)$ on a MAG $H$, we get a PAG that only has circles on
some $I_d \tuh d$ edges. For each orientation of that PAG, we can $\doit(B)$ (with $B, D$ disjoint). Here the
ordering doesn't matter, we get $H_{\doit(B),\bar{D}}$.}

We want to show:
\begin{Con}\label{lem:zhang_2008_lemma_12_alt}
  Let $H$ be a MAG with input nodes $J$ and output nodes $V$, that represents some CADMG with input nodes $J$ and output nodes $V^+ = V \dcup S$ given $S$.
  Let $B, D \subseteq V$ be disjoint.
  Define $G^* := G_{\doit(I_B,D)}$ and $H^* := \MAG(G^* \given S)$.
  For $a, b \in V \cup J$ and $C \subseteq V \cup J$:
  If there is an $m$-connecting path between $a$ and $b$ given $C$ in $H^*$ then there is one in $H_{\doit(I_D),\bar{D}}$.
\end{Con}

\subsubsection{Old stuff (should cleanup!)}

Let's consider DMAGs for now.
\cite{Borboudakis++2012} note that for DMAGs, each directed edge in a DMAG is either definitely direct (i.e., a directed edge must be present in each underlying ADMG) or possibly-indirect (i.e., there exists an underlying in ADMG in which the directed edge is absent), and give equivalent conditions for both cases.
\begin{Prp}\label{prp:dmag_possibly-indirect}
  Let $M$ be a DMAG and $i \tuh j$ a directed edge in $M$.
  \begin{enumerate}
    \item If $i \tuh j$ is invisible: it is possibly-indirect iff there exists a directed path from $i$ to $j$ in $M$ that avoids the edge $i \tuh j$.
    \item If $i \tuh j$ is visible: it is possibly-indirect iff there exists $k$ in $M$ such that $i \tuh k \tuh j$ in $M$ with $i \tuh k$ visible and $k \tuh j$ invisible.
  \end{enumerate}
\end{Prp}
\begin{proof}
  \begin{enumerate}
    \item Suppose $i \tuh j$ in $M$ is invisible.

      ``$\implies$'': If $i \tuh j$ is possibly-indirect, then $i \tuh j$ is absent in some $G$ represented by $H$. 
      Yet, we need that $i \in \Anc_G(j)$.
      Therefore there must be some directed path from $i$ to $j$ in $G$ that avoids the edge $i \tuh j$.
      That same directed path must be present in $M$.

      ``$\impliedby$'': Suppose there exists a directed path from $i$ to $j$ in $M$ that avoids the edge $i \tuh j$.
      We can take $G'$ obtained from $M$ by removing $i \tuh j$ and adding $i \huh j$.
      Then $i \in \Anc(j)$ in $G'$ and there is an inducing path between $i$ and $j$ in $G'$.
      One can check that $G'$ is represented by $H$. \Joris{TODO: this needs work! It's an extension of Lemma~\ref{lem:zhang2008_lemma_10_mag_extended}.}
      Since $G'$ doesn't contain $i \tuh j$, the edge $i \tuh j$ is possibly-indirect.

  \item Suppose $i \tuh j$ in $M$ is visible.

      ``$\implies$'': If $i \tuh j$ is possibly-indirect, then $i \tuh j$ is absent in some $G$ represented by $H$. 
      There must be an inducing path between $i$ and $j$ in $G$ into $j$.
      Because $i \tuh j$ is visible, no inducing path in $G$ between $i$ and $j$ can be into $i$.
      Hence there is an inducing path in $M$ out of $i$ into $j$.
      It must be an inducing path of the form $i \tuh k \huh \dots \huh j$, since $i \tuh j$ is not present in $G$.
      Since $i \in \Anc_G(j)$, and $k \in \Anc_G(\{i,j\})$, $k \in \Anc_G(j)$. 
      Therefore, the subpath $k \huh \dots \huh j$ must also be inducing, and $k \tuh j$ is present in $M$. 
      Since the inducing path $k \huh \dots \huh j$ is into $k$, $k \tuh j$ must be invisible in $M$.
      Furthermore $i \tuh k$ is present in $M$.
      If it were invisible, we cannot rule out the existence of $i \huh k$ in $G$.
      But that would give an inducing path $i \huh k \huh \dots \huh j$ into $i$, contradicting the visibility of $i \tuh j$.
      Hence, $i \tuh k$ must be visible in $M$.

      So: for visible $i \tuh j$ we have shown that if it is possibly-indirect, i.e., there exists $G$ represented by $H$ without $i \tuh j$, then there must exist $k$ in $M$ such that $i \tuh k \tuh j$ in $M$ and $k \tuh j$ is invisible and $i \tuh k$ is visible.

      ``$\impliedby$'': Under the assumptions, there exists a $G$ represented by $H$ with $k \huh j$ present.
      If we remove the edge $i \tuh j$ from $G$ (if present), then the resulting $G'$ still has an inducing path between $i$ and $j$, and still has $i \in \Anc(j)$, and therefore still gives $i \tuh j$ in its representative MAG.
      Note that if we take $G'$ obtained from $H$ by removing $i \tuh j$ and adding $k \huh j$ it fulfills all requirements ($G'$ is represented by $H$, $G'$ does not contain $i \tuh j$).
      Hence $i \tuh j$ is possibly-indirect.
  \end{enumerate}
\end{proof}
See also Figure~\ref{fig:mag_do} and Figure~\ref{fig:mag_do2}.
What holds for each $H^*$?
\Joris{(Below, take $S = \emptyset$!)}
\begin{Prp}
  Let $G$ be a CADMG with input nodes $J$ and output nodes $V^+ = V \dcup S$. 
  Denote the DMAG that represents $G$ given $S$ by $H := \MAG(G \given S)$.
  For $D \subseteq V$, define $G^* := G_{\doit(D)}$ and $H^* := \MAG(G_{\doit(D)} \given S)$.
  Then:
  \begin{enumerate}
    \item the skeleton of $H^*$ is a subgraph of the skeleton of $H_{\bar D}$;
      (hence $H^*$ has no edges into $D$);
    \item for $i,j$ adjacent in $H^*$:
      \begin{itemize}
        \item if $i \suh j$ in $H$ then $i \suh j$ in $H^*$
        \item if $i \huh j$ in $H$ then $i \huh j$ in $H^*$
        \item if $i \tuh j$ in $H$ is definitely direct, then $i \tuh j$ in $H^*$
        \item if $i \tuh j$ in $H$ is visible, then $i \tuh j$ in $H^*$
      \end{itemize}
    \item if $i \tus j$ in $H^*$ then $i \tus j$ in $H$;
  \end{enumerate}
%  Then: for $i,j \in V \cup J$: if $i,j$ are adjacent in $H^*$ then $i,j$ are adjacent in $H_{\bar D}$,
%  and the edge in $H^*$ is into $i$ iff the edge in $H_{\bar D}$ is into $i$, 
%  and the edge in $H^*$ is into $j$ iff the edge in $H_{\bar D}$ is into $j$. 
\end{Prp}
\begin{proof}
  Note that $x \in \Anc_{G_{\doit(D)}}(U)$ implies that $x \in \Anc_G(U)$, for all $x \in V^+ \cup J$ and $U \subseteq V^+ \cup J$.

  \begin{enumerate}
    \item
      Let $i \sus j$ be an edge in $H^*$. 
      Then there is an inducing path $\pi$ between $i$ and $j$ given $S$ in $G_{\doit(D)}$.
      Every non-endpoint node on the path must be a collider in $\Anc_{G_{\doit(D)}}(\{i,j\} \cup S)$.
      The same sequence of node indices defines a path in $G$ (with possibly different node types and/or edges), where every non-endpoint node must be a collider in $\Anc_G(\{i,j\} \cup S)$.
      Therefore, there is an inducing path between $i$ and $j$ given $S$ in $G$.
      Hence, $i$ and $j$ are adjacent in $H$.

      (If $H^*$ had an edge $i \suh d$ for $d \in D$ then in $G^*$ there must be an inducing path between $i$ and $d$ given $S$ that is \emph{into} $d$.
      However, $G^*$ contains no edges into $d$.
      Contradiction.
      So $H^*$ has no edges into $D$.)
    \item
      \begin{itemize}
        \item
          Any arrowhead $i \suh j$ on an edge in $H$ (sic!) signifies that
          $j \notin \Anc_G(\{i\} \cup S)$, hence $j \notin \Anc_{G_{\doit(D)}}(\{i\} \cup S)$, and hence $i \suh j$ in $H^*$ if $i,j$ are adjacent in $H^*$.
        \item If $i \tuh j$ in $H$ is definitely direct, it must be present in all $G$ represented by $H$.
          If additionally $j \notin D$, it must be present in all $G_{\doit(D)}$ where $G$ is represented by $H$.
          Hence it must be present in all possible $H^*$, and therefore we can orient $i \tuh j$ in $H_{\doit(D)}$.
        \item Suppose $i \tuh j$ in $H$. 
          If it is visible, then $i \tuh j$ in $H^*$.
          Proof by contradiction.
          An edge $i \suh j$ must be in $H^*$.
          Suppose it is $i \huh j$.
          Then there is an inducing path in $G_{\doit(D)}$ between $i$ and $j$ given $S$ that is into $i$ and into $j$.
          This implies that $i,j \notin D$ and moreover that there is an inducing path in $G$ between $i$ and $j$ given $S$ that is into $i$ and into $j$.
          %The edge $i \sus j$ in $H$ is either $i \huh j$, or (by Lemma~\ref{lem:zhang2008_lemma_9_mag_extended}) an \emph{invisible} directed edge $i \tuh j$ or $i \hut j$.
          This gives a contradiction with Lemma~\ref{lem:zhang2008_lemma_9_mag_extended}.
          So the edge in $H^*$ must be $i \tuh j$.
      \end{itemize}
    \item
      Suppose the edge $i \sus j$ in $H^*$ were out of $i$.
      Then $i \in \Anc_{G_{\doit(D)}}(\{j\} \cup S)$, and hence $i \in \Anc_G(\{j\} \cup S)$.
      Therefore, the edge $i \sus j$ in $H$ must be out of $i$ as well.
      Similarly, if the edge $i \sus j$ in $H^*$ were out of $j$, the edge $i \sus j$ in $H$ must be out of $j$ as well.
  \end{enumerate}
\end{proof}
\Joris{The remaining case:
  Suppose $i \tuh j$ in $H$ is possibly-indirect and invisible, and there are directed paths from $i$ to $j$ in $H_{\bar D}$ except for $i \tuh j$.
  Consider all directed paths from $i$ to $j$ in $H_{\bar D}$ except for $i \tuh j$.
  If one of them is already oriented as a directed path in $H_{\doit(D)}$, then orient $i \tuh j$.
  Otherwise, all of them contain at least one unoriented edge or one oriented-as-possibly-bidirected edge in $H_{\doit(D)}$.
  If all of them contain one oriented-as-possibly-bidirected edge in $H_{\doit(D)}$ then OPTIMISTICALLY orient $i \ouh j$ as well.
  The remaining case is that some of them contain an unoriented edge.
  Pick one, say $a \tuh b$.
  It must be possibly-indirect and invisible, and there are directed paths from $a$ to $b$ in $H_{\bar D}$ expect for $a \tuh b$.
  We can now recurse into orienting that unoriented edge.
  Then proceed to the next unoriented edge on the directed path, and recurse into orienting it, and repeat for all unoriented edges on the directed path.
  Finally, all edges on the directed path that were unoriented, have been oriented.
  Then proceed in a similar way with another partially-unoriented directed path, until they all have been oriented.
  We can then orient the edge $i \tuh j$ accordingly.

  Note that at some point deep in the recursion, we have run out of nodes to substitute an unoriented edge by a directed path in $H_{\bar D}$.
  Indeed, because of the acyclicity of $H_{\bar{D}}$, the directed paths we substitute for an unoriented edge need to go via descendants of a, and ancestors of b; so these nodes cannot overlap with either ancestors of a or descendents of b.
  Therefore, the ``walkaround'' path goes via nodes that were not already on the directed path.
  At some point, we simply run out of available nodes.
  That means that recursion stops there, and that unoriented edge will get oriented.

  Now what we have to check is whether if we replace in the final $H_{\doit(D)}$ all possibly-bidirected edges by bidirected ones, that we can find a $G$ that represents it (itself?). 
  Or much stronger even: for any subset of possibly-bidirected edges, there exists $G$ represented by $H$ such that the corresponding $H^*$ has bidirected edges precisely on this subset, and otherwise, directed edges.
  In that case, we can consider $H_{\doit(D)}$ a superPAG.
  The first property might even hold...!
  For the innermost possibly-bidirected edges, we can use the following proposition.
  This means that if we remove corresponding \emph{directed} edges (make them surely-indirect) then such a representative ADMG exists.
  Tracing back in the recursion, we then oriented edges for which all ``walkaround'' paths contained at least one possibly-bidirected edge, so we know that none of them survives in $G_{\doit(D)}$ (by a similar argument as in the proposition below).
  However, what we still have to show is that if we have two possibly-indirect edges, then there exists a common representative ADMG in which both are indeed indirect.
  This seems to work, but I'm not really sure right now.
  We could then selectively add direct $i \tuh j$ edges to change an $H^*$ into another $H^*$ by replacing one bidirected by one directed edge.
  If we can do so, it suggests all orientations of $H_{\doit(D)}$ are actually valid MAGs.
  If we can then show that they are Markov equivalent, we are done.

  This looks good!
  Note that we can also formulate the existence result without defining the recursive algorithm:
  As long as there exists an unoriented $i \tuh j$ in $H$ which is possibly-indirect and invisible, and there are directed paths from $i$ to $j$ in $H_{\bar D}$ except for $i \tuh j$,
  then there exists such an unoriented $i \tuh j$ in $H$ which is possibly-indirect and invisible and for which either
  (i) there exists a directed path from $i'$ to $j'$ in $H_{\bar D}$ except for $i \tuh j$ that is already oriented as directed path, or
  (ii) all directed paths from $i'$ to $j'$ in $H_{\bar D}$ except for $i \tuh j$ contain at least one oriented-as-possibly-bidirected edge in $H_{\doit(D)}$.
  In the former case, we can orient $i' \tuh j'$ in $H_{\doit(D)}$, in the latter case $i' \ouh j'$ in $H_{\doit(D)}$.
  We now show the existence of such $i' \tuh j'$.
  Write $i_1 := i$, $j_1 = j$.
  If $i_1 \tuh j_1$ doesn't qualify, then there must be a directed path $\pi_1$ from $i_1 \tuh j_1$ (avoiding the edge $i_1 \tuh j_1$ and $D$) that has at least one unoriented edge $i_2 \tuh j_2$.
  The non-endpoint nodes on that directed path cannot overlap with $\Anc_H(i_1) \cup \Desc_H(j_1)$ (otherwise $H$ would contain a cycle).
  Hence the path avoids directed edges into $D \cup \Anc_H(i_1) \cup \Desc_H(j_1)$.
  Note that the path $\pi_1$ consists at least of 3 nodes (including $i_1, j_1$).
  Now consider the unoriented edge $i_2 \tuh j_2$.
  Either it qualifies, or otherwise, there must be a directed path from $i_2 \tuh j_2$ (avoiding the edge $i_2 \tuh j_2$ and $D \cup \Anc_H(i_1) \cup \Desc_H(i_2)$) that has at least one unoriented edge $i_3 \tuh j_3$.
  That path avoids directed edges into $D \cup \Anc_H(i_1) \cup \Desc_H(j_1) \cup \Anc_H(i_2) \cup \Desc_H(j_2)$, otherwise $H$ would contain a cycle.
  Note that the path obtained from $\pi_1$ by substituting the edge $i_2 \tuh j_2$ by the path $\pi_2$ consists at least of 4 nodes.
  We can repeat this procedure, until either we find a qualifying edge $i_n \tuh j_n$, or we run out of nodes.

  Let us keep track of the order in which we orient possibly-indirect and invisible edges in $H$ for which there are directed paths from $i$ to $j$ in $H_{\bar D}$ that avoid $i \tuh j$ itself.
  Every time we decide to orient such an edge as $i \ouh j$ in $H_{\doit(D)}$, there exists a walkaround path from $i$ to $j$, but none that is already oriented as a directed path in $H_{\doit(D)}$, and each such path contains another edge oriented already as $i' \ouh j'$ in $H_{\doit(D)}$.
  We know that if we replace $i' \ouh j'$ by $i' \huh j'$ there is still a directed path from $i'$ to $j'$ in $H$ that avoids $i' \tuh j'$ (but not in $H_{\doit(D)}$).
  For the first orientation as possibly-bidirected, we know that when replacing this edge in $H$ by a bidirected one, there is still a corresponding indirect directed path in $H$ left.
  But what if one of the edges on that path later becomes oriented as possibly-bidirected?
  So perhaps we should only do the orientation as possibly-bidirected if we can guarantee that there is still a corresponding indirect directed path in $H$ left that will never disappear by subsequent orientations?
  In other words, only orient $i \tuh j$ as possibly-bidirected if there is an indirect directed path in $H$ for which each edge is either (i) into $D$, or (ii) already oriented as directed in $H_{\doit(D)}$.
  Or, perhaps, we were too optimistic.
  Perhaps one can show that by turning some of the innermost edges into bidirected, some later on need to be directed in order for there to remain a directed path.

  Hm.
  Perhaps the following orientation algorithm works itself inside-out in such a way that it works.
  When the procedure $\mathtt{Orient}(i \tuh j)$ is called on a possibly-indirect invisible edge, then it will first invoke $\mathtt{Orient}(k \tuh l)$ for each yet unoriented edge on any indirect directed path from $i$ to $j$ in $H$.
  The innermost orientation steps are to orient definitely-direct edges (as directed), and visible edges (as directed).
  These do remain present in $G'$.
  What remains to orient are the possibly-indirect invisible edges in $H$.
  From those, we should first orient the ones as directed.
  Aha, we could do that recursively as a first pass, and then we have obtained all directed edges in $H_{\doit(D)}$.
  What remains must be the possibly-bidirected edges in $H_{\doit(D)}$.
  Really???
  Now we can also tackle them ``inside-out''.
  So first consider the ones that don't have a walkaround path in $H$ that contains another possibly-bidirected edge.
  These only have walkaround paths that consist of edges that are oriented as directed in $H_{\doit(D)}$ and therefore these paths are present in $G'$ (yes!).
  Mark these possibly-bidirected edges as ``done''.
  Then proceed with the remaining possibly-bidirected edges.
  Some of them will have a walkaround path that consists only of certainly-directed edges and ``done'' possibly-bidirected edges.
  Then we can mark these possibly-bidirected edges as ``done'' as well.
  In this way we can proceed until all possibly-bidirected edges are marked as ``done''.
  We have shown that if we replace each possibly-bidirected edge by a bidirected one, we still preserve all ancestral relations of $H$.
  Idea: if we use this ordering to recursively orient edges, we squat two flies with a single hit.

  Let's try to prove completeness of this definition of $H_{\doit(D)}$.

  First we unveil some structure behind the ``invisible, possibly indirected'' directed edges in a MAG.
  \begin{Def}
    Let $G$ be a CADMG with input nodes $J$ and output nodes $V^+ = V \dcup S$ and $H$ the MAG that represents $G$ given $S$.
    For now we assume $S = \emptyset$, so $H$ is a DMAG.
    We define a directed graph $\C{K}(H)$ as follows:
    \begin{itemize}
      \item The nodes of $\C{K}$ are the invisible and possibly-indirect directed edges in $H$.
      \item For any $i \tuh j$ in $\C{K}$: consider the set of all directed paths in $H$ from $i$ to $j$ that avoid $i \tuh j$, and then for each edge $k \tuh l$ on such a path that is in $\C{K}$, declare it to be a parent of $i \tuh j$.
    \end{itemize}
    This defines a DAG because if there were a cycle in this DAG, there would be a cycle in the MAG as well.
  \end{Def}
  By traversing the edges in $\C{I} \subseteq \C{K}$ according to that topological ordering, we can show that every time we replace an edge in $\C{I}$ with a bidirected one, the corresponding ancestral relation is still guaranteed to be preserved. 
  So in this way we can show that \emph{any} subset of replacements $\C{I} \subseteq \C{K}$ (of a directed by a bidirected edge in $H$) preserves the ancestral relations of $H$. 
  \begin{Prp}
    Let $\C{I} \subseteq \C{K}$ be a subset of invisible and possibly-indirect directed edges in $H$.
    Order $\C{I}$ according to a topological ordering of $\C{K}$: $i_1 \tuh j_1, i_2 \tuh j_2, \dots, i_n \tuh j_n$ (from root nodes to sink nodes in the DAG $\C{K}$).
    Define a sequence of $n$ ADMGs:
    $G_0 := H$, $G_t$ obtained from $G_{t-1}$ by replacing $i_t \tuh j_t$ by $i_t \huh j_t$.
    Each ADMG in this sequence has exactly the same (non-)ancestral relations as $H$.
  \end{Prp}
  \begin{proof}
    We will show that the ancestral relations of $G_t$ coincide with that of $G_{t-1}$ for $t=1,\dots,n$.

    Let us start with $t=1$.
    $G_1$ is obtained by replacing in $G_0$ $i_1 \tuh j_1$ by $i_1 \huh j_1$.
    As there is a directed path from $i_1$ to $j_1$ in $G_0$ that avoids the edge $i_1 \tuh j_1$, the ancestral relations do not change by the replacement.
    Note that this directed path will not be removed in a future replacement step (as it contains no edges in $\C{I}$), so this directed path will remain present in all $G_r$.

    Now consider any $t > 1$.
    We know that there is a directed path from $i_t$ to $j_t$ in $\C{K}$ that avoids $i_t \tuh j_t$.
    It may contain directed edges that are in $\C{I}$.
    These will have been replaced by bidirected edges in earlier stages.
    We can replace those with permanent directed paths in $G_{t-1}$.
    Consider an ancestor $k \tuh l$ of $i_t \tuh j_t$ in $\C{K}$.
    Either it is not in $\C{I}$, and then the directed edge $k \tuh l$ is present in $G_{t-1}$, or it is in $\C{I}$, and then we can replace it by a permanent directed path.
    This gives us a (permanent) directed path from $i_t$ to $j_t$ that will remain present in all $G_r$.
    So we can safely replace $i_t \tuh j_t$ by $i_t \huh j_t$.
  \end{proof}
  For the specific $\C{I}$ we want to use (all $i \ouh j$ edges in $H_{\doit(D)}$) we can show in addition that their replacement leads to $i$ no longer being an ancestor of $j$ in $G_{\doit(D)}$. So we get a representative of a maximally-bidirected orientation of $H_{\doit(D)}$ and one of a maximally-directed orientation (just add back all directed edges that were replaced by bidirected edges).}

\Joris{
\begin{Lem}\label{lem:zhang2008_lemma_10_mag_extended_further}
  Let $G$ be a CADMG with input nodes $J$ and output nodes $V^+ = V \dcup S$ and $H$ the MAG that represents $G$ given $S$.
  %Let $Z \subseteq V \cup J$ such that each of which has no incoming arrowhead in $H$ (one can always take $Z = \emptyset$).
  %Let $i \tuh j$ be a directed edge in $H$.
  Let $\C{I} = \{i \tuh j\}$ be a set of invisible possibly-indirect directed edges in $H$.\footnote{I think we need to assume in addition that $\C{I}$ has the property that if we replace in $H$ all $i \tuh j$ in $\C{I}$ by $i \huh j$, all ancestral relations in $H$ are preserved. Ah, it seems I have just shown that this property holds automatically for any such subset $\C{I}$!}
  That is, each $i \tuh j$ is in $H$, is invisible, and $H$ contains a directed path $\pi_{i \tuh j}$ from $i$ to $j$ that avoids the edge $i \tuh j$.
  Then there exists a CADMG $\tilde{G}$ with input nodes $J$ and output nodes $\tilde{V}^+ = V \dcup \tilde{S}$ such that:
  \begin{enumerate}
    \item $\tilde{G}$ is also represented by $H$ given $\tilde{S}$, and
    \item for each $i \tuh j \in \C{I}$, $i \huh j$ is in $\tilde{G}$ and $i \tuh j$ is not in $\tilde{G}$.
  \end{enumerate}
\end{Lem}
\begin{proof}
  We construct CADMG $\tilde{G}$ as follows.
  It has nodes $\tilde{V}^+ \cup J$ with $\tilde{V}^+ = V \cup \tilde{S}$, where $\tilde{S} := \{s_{\{x,y\}} : x \tut y \in H\}$.
  We add all bidirected edges $x \huh y$ for $x,y \in V\cup J$ that are present in $H$ to $\tilde{G}$.
  In addition, we add for each $i \tuh j \in \C{I}$ the bidirected edge $i \huh j$ to $\tilde{G}$.
  We then add all directed edges $x \tuh y$ for $x,y \in V\cup J$ that are present in $H$, except for those in $\C{I}$, to $\tilde{G}$.
  For the undirected edges $x \tut y$ in $H$, we add $x \tuh s_{\{x,y\}} \hut y$ to $\tilde{G}$.
  Note that $\tilde{G}$ is acyclic by construction.

  Since all directed edges in $H$ are also present in $\tilde{G}$, and no directed edge absent in $H$ is present in $\tilde{G}$, the ancestral relations between variables in $V \cup J$ are the same in $H$ as in $\tilde{G}$. 
  \Joris{This is a bit more tricky now. We need to stick to an ordering for the edges in $\C{I}$.
  Indeed, if we can arrange that $\C{I}$ is ordered such that $i \tuh j < k \tuh l$ iff $\pi_{i \tuh j}$ avoids $k \tuh l$ (???) we can show that ancestral relations are preserved in $\tilde{G}$.
  It is trivial that each directed edge present in $\tilde{G}$ is present in $H$.}
  That is, for $x,y \in V \cup J$, $x \in \Anc_H(y) \iff x \in \Anc_{\tilde{G}}(y)$.
  
  We now show that for $v \in V \cup J$: if $v \in \Anc_{\tilde{G}}(\tilde{S})$, then $v \in \Anc_G(S)$ (note: the converse does not hold).
  Indeed, if $v \in \Anc_{\tilde{G}}(\tilde{S})$ then there must be a directed path $v \tuh \dots \tuh w \tuh s$ in $\tilde{G}$ such that $s \in \tilde{S}$, and all other nodes on the directed path are in $V \cup J$.
  The directed path $v \tuh \dots \tuh w$ must also be present in $H$. 
  The node $s$ must be $s_{\{w,u\}}$ for some $u \in V \cup J$ with $w \tut u$ in $H$.
  But since $\dots \suh w \tut u$ cannot occur in a MAG $H$ (because $G$ is acyclic), the directed path must be the trivial path $v = w$.
  Hence $v \tut u$ in $H$, and therefore $v \in \Anc_G(S)$.
  We also state the contra-positive for convenience: for $v \in V \cup J$: if $v \notin \Anc_G(S)$ then $v \notin \Anc_{\tilde{G}}(\tilde{S})$.
  \Joris{Furthermore, observe that if $x \hus y$ in $H$ then $x \notin \Anc_G(\{y\} \cup S)$ and hence $x \notin \Anc_{\tilde{G}}(\{y\} \cup \tilde{S})$.}

  It remains to show that $H$ represents $\tilde{G}$ given $\tilde{S}$.
  For that we need to show:
  \begin{enumerate}
    \item For all $x \in V \cup J, y \in V$: $x \sus y$ is in $H$ if and only if there is an inducing path between $x$ and $y$ given $\tilde{S}$ in $\tilde{G}$;
    \item If $x \hus y$ in $H$ then $x \notin \Anc_{\tilde{G}}(\{y\} \cup \tilde{S})$;
    \item If $x \tus y$ in $H$ then $x \in \Anc_{\tilde{G}}(\{y\} \cup \tilde{S})$;
  \end{enumerate}

  First adjacency.
  Let $x \in V \cup J$, $y \in V$.
  Suppose $x$ and $y$ are not adjacent in $H$.
  We must show that in $\tilde{G}$ there is no inducing path between $x$ and $y$ given $\tilde{S}$.
  Suppose there is one.
  Pick one of minimal length.
  It must be a collider path in which every collider is in $\Anc_{\tilde{G}}(\{x,y\} \cup \tilde{S})$.

  Suppose the path contains a node of $\tilde{S}$.
  Then the path has to be of the form $x \dots a \tuh s_{\{a,b\}} \hut b \dots y$ with $a \tut b$ in $H$.
  Since only the endpoint nodes can be non-colliders, the path must be $x \tuh s_{\{x,y\}} \hut y$ with $x \tut y$ in $H$. 
  Contradiction.
  We conclude that the path can only contain nodes in $V \cup J$.

  Assume that each edge on the path is present in $H$ as well (in other words, no bidirected edge corresponding to an edge in $\C{I}$).
  Then each collider on the path is not in $\Anc_G(S)$.
  Therefore, by the result above, each collider on the path is not in $\Anc_{\tilde{G}}(\tilde{S})$.
  But this means that each collider on the path must be in $\Anc_{\tilde{G}}(\{x,y\})$, and hence in $\Anc_H(\{x,y\})$.
  So then the path is also an inducing path between $x$ and $y$ in $H$, which contradicts the maximality of $H$.

  Now suppose some $i \huh j$ obtained by replacing an edge in $\C{I}$ is on the path. 
  Pick the first such edge on the path encountered from $x$.
  %All other edges on the path must also be present in $H$.
  The edge $i \huh j$ is not in $H$, but the edge $i \tuh j$ or $i \hut j$ or $i \tut j$ is in $H$.
  %This corresponds with the directed edge $i \tuh j$ in $H$.
  In the latter case, there must be a path of the form $x \suh \dots \huh \dots \huh i \tut j \huh \dots \huh \dots \hus y$ in $H$. (This needs fixing!)
  But this can only happen for $x=i$ and $y=j$, that is, if the path consists of the single edge $i \tut j$.
  But then $x$ would be adjacent to $y$ in $H$, a contradiction.

  Hence the only possibilities are that $i \tuh j$ in $H$ or $i \hut j$ in $H$.
  Let us first consider the case that the inducing path is of the form $x \suh \dots \huh \dots \huh i \huh j \dots y$ with all edges between $x$ and $i$ present in $H$,
  and that $i \tuh j$ is in $H$.
  If $x = i$ then we can use the same reasoning as before to show that the path $x = i \tuh j \dots y$ in $H$ must be inducing, contradicting the maximality of $H$. UPDATE: THIS MAY FAIL NOW!
  THE PROBLEM SEEMS TO BE THAT ALL THE BIDIRECTED EDGES INTRODUCED IN $\tilde{G}$ MAY LEAD TO ADDITIONAL INDUCING PATHS THAT ARE NOT IN $H$.
  Well, the path $x = i \tuh j \dots y$ is still inducing, but it is not in $H$.
  We could even try to show that each of these $G_i$'s is itself a MAG (representing itself).
  If $x \ne i$, then the first part of the inducing path in $\tilde{G}$, say $z_0 \dots z_k$ (with $z_0 = x$ and $z_k = i$ and $k \ge 1$) is also a collider path into $i$ in $H$.
  We will prove by induction that for $1 \le s \le k-1$, $z_s$ is a parent of $j$ in $H$.

  AN IDEA TO MITIGATE THE PROBLEM: WE SHOW THAT THE SEQUENCE $G_0, \dots, G_n$ of ADMGs are represented by the same skeleton, by relating the existence of inducing paths between them.
  A single replacement is OK: we can show that if $x$ and $y$ are not adjacent in $H = G_0$, then in $G_1$ there is no inducing path between $x$ and $y$ given $\tilde{S}$.
  Note that $x$ and $y$ in $J \cup V$ are adjacent in $H$ if and only if they are adjacent in $G_r$ for any $r = 0, \dots, n$.
  Can we show that if $x$ and $y$ are not adjacent in $G_1$, then in $G_2$ there is no inducing path between $x$ and $y$ given $\tilde{S}$?
  If we can, we can probably show that if $x$ and $y$ are not adjacent in $G_r$, then in $G_{r+1}$ there is no inducing path between $x$ and $y$ given $\tilde{S}$.
  \textbf{Generally: can we show that if there is no inducing path between $x$ and $y$ in $G_r$, then there is no inducing path between $x$ and $y$ in $G_{r+1}$?}
  It would be convenient if the following properties are preserved in the sequence: the edges in $\C{I}$ are invisible in each $G_r$ (this is suggested by Zhang
  in the proof of Lemma 12 in 2008), and each $G_r$ is ancestral.
  So in $H$ we have: $i \tuh j$ is invisible, $a \tuh b$ is invisible, and now we change $i \tuh j$ into $i \huh j$, could that make $a \tuh b$ visible?

  Now suppose that for some $0 \le s \le k-1$, $z_{s+1} \tuh j$ is in $H$ (note this holds in particular for $s=k-1$). \Joris{No, we need to assume $z_t \tuh j$ is in $H$ for all $t = s+1, \dots, k$ to get a valid induction.} 
  We will show that $z_s \tuh j$ is in $H$.
  $z_s$ must be adjacent to $j$ in $H$, otherwise $i \tuh j$ would be visible in $H$.
  The edge between $z_s$ and $j$ cannot be $z_s \hut j$, otherwise $j \tuh z_s \suh z_{s+1} \tuh j$ in $H$, and $H$ would be non-ancestral.  
  It can also not be $z_s \huh j$, otherwise there would be an inducing path $x = z_0 \dots z_s \huh j \dots y$ in $\tilde{G}$ that contains fewer nodes, a contradiction.
  It can also not be $z_s \tut j$, as that would give $i \tuh j \tut z_s$ in $H$ which is not possible.
  Hence it has to be $z_s \tuh j$.

  So by induction, for $1 \le s \le k-1$, $z_s$ is a parent of $j$ in $H$.
  So this gives a collider path $x \suh z_1 \huh \dots \huh z_{k-1} \huh i$ in $H$ such that each $z_s$ is a parent of $j$ in $H$, and $i \tuh j$ in $H$.
  Therefore, $x$ must be adjacent to $j$ in $H$, otherwise $i \tuh j$ would be visible.
  If this edge were of the form $x \hut j$ then $H$ would not be ancestral (because $z_1 \tuh j \tuh x \suh z_1$ in $H$).
  The edge cannot be of the form $x \tut j$ either because then $z_1 \tuh j \tut x$ would be in $H$.
  Hence it must be of the form $x \suh j$.
  But then this edge must also be present in $\tilde{G}$.
  Then we have an inducing path $x \suh j \dots y$ in $\tilde{G}$ that is shorter than the one that was assumed to be of minimal length.
  Contradiction.

  We have shown that if $x$ and $y$ are not adjacent in $H$, there is no inducing path between $x$ and $y$ in $\tilde{G}$ given $\tilde{S}$.
  Now suppose $x$ and $y$ are adjacent in $H$.
  We distinguish four cases corresponding to the type of the edge between $x$ and $y$ in $H$. 
  \begin{itemize}
    \item $x \tuh y$: By construction $x \tuh y$ is in $\tilde{G}$. \Joris{optionally: for $x \tuh y$ equals $i \tuh j$, we have $i \huh j$ in $\tilde{G}$ instead}
      Then $y \notin \Anc_G(S)$, and therefore also $y \notin \Anc_{\tilde{G}}(\tilde{S})$.
      Therefore, $x \in \Anc_{\tilde{G}}(\{y\} \cup \tilde{S})$ and $y \notin \Anc_{\tilde{G}}(\{y\} \cup \tilde{S})$.
    \item $x \hut y$: This case is similar (with the roles of $x,y$ swapped).
    \item $x \huh y$: By construction $x \huh y$ is in $\tilde{G}$.
      Then both $x, y \notin \Anc_G(S)$, and therefore also $x,y \notin \Anc_{\tilde{G}}(\tilde{S})$. 
      Therefore, $x \notin \Anc_{\tilde{G}}(\{y\} \cup \tilde{S})$ and $y \notin \Anc_{\tilde{G}}(\{y\} \cup \tilde{S})$.
    \item $x \tut y$: Then $x \tuh s_{\{x,y\}} \hut y$ in $\tilde{G}$.
      Therefore, $x \in \Anc_{\tilde{G}}(\{y\} \cup \tilde{S})$ and $y \in \Anc_{\tilde{G}}(\{y\} \cup \tilde{S})$.
  \end{itemize}
  In all cases, there is an inducing path between $x$ and $y$ given $\tilde{S}$ in $\tilde{G}$.
  We conclude that $H$ represents $\tilde{G}$ given $\tilde{S}$.
\end{proof}
}

\begin{figure}\centering
  \begin{tikzpicture}
    \begin{scope}[xshift=-4cm]
      \node[ndout] (a) at (0,0) {$a$};
      \node[ndout] (b) at (0,-1.5) {$b$};
      \node[ndout] (z) at (-1.5,0) {$z$};
      \node[ndout] (d1) at (-0.75,1) {$d_1$};
      \node[ndout] (d2) at (1,-0.75) {$d_2$};
      \draw[tuh] (a) edge (z);
      \draw[tuh] (a) edge (d1);
      \draw[tuh] (d1) edge (z);
      \draw[tuh] (a) edge (b);
      \draw[tuh] (a) edge (d2);
      \draw[tuh] (d2) edge (b);
      \node at (0,-2.5) {$H$};
    \end{scope}
    \begin{scope}
      \node[ndout] (a) at (0,0) {$a$};
      \node[ndout] (b) at (0,-1.5) {$b$};
      \node[ndout] (z) at (-1.5,0) {$z$};
      \node[ndout] (d1) at (-0.75,1) {$d_1$};
      \node[ndout] (d2) at (1,-0.75) {$d_2$};
      \draw[tuh] (a) edge (z);
      \draw[tuh] (a) edge (d1);
      \draw[tuh] (d1) edge (z);
      \draw[tuh] (a) edge (b);
      \draw[tuh] (a) edge (d2);
      \draw[tuh] (d2) edge (b);
      \node at (0,-2.5) {$G$};
    \end{scope}
    \begin{scope}[xshift=4cm]
      \node[ndout] (a) at (0,0) {$a$};
      \node[ndout] (b) at (0,-1.5) {$b$};
      \node[ndout] (z) at (-1.5,0) {$z$};
      \node[ndint] (d1) at (-0.75,1) {$d_1$};
      \node[ndint] (d2) at (1,-0.75) {$d_2$};
      \draw[tuh] (a) edge (z);
      \draw[tuh] (d1) edge (z);
      \draw[tuh] (a) edge (b);
      \draw[tuh] (d2) edge (b);
      \node at (0,-2.5) {$G_{\doit(X)}$};
    \end{scope}
    \begin{scope}[xshift=8cm]
      \node[ndout] (a) at (0,0) {$a$};
      \node[ndout] (b) at (0,-1.5) {$b$};
      \node[ndout] (z) at (-1.5,0) {$z$};
      \node[ndint] (d1) at (-0.75,1) {$d_1$};
      \node[ndint] (d2) at (1,-0.75) {$d_2$};
      \draw[tuh] (a) edge (z);
      \draw[tuh] (d1) edge (z);
      \draw[tuh] (a) edge (b);
      \draw[tuh] (d2) edge (b);
      \node at (0,-2.5) {$\MAG(G_{\doit(X)}) = H^*$};
    \end{scope}
    \begin{scope}[yshift=-4cm]
      \node[ndout] (a) at (0,0) {$a$};
      \node[ndout] (b) at (0,-1.5) {$b$};
      \node[ndout] (z) at (-1.5,0) {$z$};
      \node[ndout] (d1) at (-0.75,1) {$d_1$};
      \node[ndout] (d2) at (1,-0.75) {$d_2$};
      \draw[huh] (a) edge (z);
      \draw[tuh] (a) edge (d1);
      \draw[tuh] (d1) edge (z);
      \draw[tuh] (a) edge (b);
      \draw[tuh] (a) edge (d2);
      \draw[tuh] (d2) edge (b);
      \node at (0,-2.5) {$G$};
    \end{scope}
    \begin{scope}[xshift=4cm,yshift=-4cm]
      \node[ndout] (a) at (0,0) {$a$};
      \node[ndout] (b) at (0,-1.5) {$b$};
      \node[ndout] (z) at (-1.5,0) {$z$};
      \node[ndint] (d1) at (-0.75,1) {$d_1$};
      \node[ndint] (d2) at (1,-0.75) {$d_2$};
      \draw[huh] (a) edge (z);
      \draw[tuh] (d1) edge (z);
      \draw[tuh] (a) edge (b);
      \draw[tuh] (d2) edge (b);
      \node at (0,-2.5) {$G_{\doit(X)}$};
    \end{scope}
    \begin{scope}[xshift=8cm,yshift=-4cm]
      \node[ndout] (a) at (0,0) {$a$};
      \node[ndout] (b) at (0,-1.5) {$b$};
      \node[ndout] (z) at (-1.5,0) {$z$};
      \node[ndint] (d1) at (-0.75,1) {$d_1$};
      \node[ndint] (d2) at (1,-0.75) {$d_2$};
      \draw[huh] (a) edge (z);
      \draw[tuh] (d1) edge (z);
      \draw[tuh] (a) edge (b);
      \draw[tuh] (d2) edge (b);
      \node at (0,-2.5) {$\MAG(G_{\doit(X)}) = H^*$};
    \end{scope}
    \begin{scope}[yshift=-8cm]
      \node[ndout] (a) at (0,0) {$a$};
      \node[ndout] (b) at (0,-1.5) {$b$};
      \node[ndout] (z) at (-1.5,0) {$z$};
      \node[ndout] (d1) at (-0.75,1) {$d_1$};
      \node[ndout] (d2) at (1,-0.75) {$d_2$};
      \draw[tuh] (a) edge (z);
      \draw[tuh] (a) edge (d1);
      \draw[tuh] (d1) edge (z);
      \draw[huh] (a) edge (b);
      \draw[tuh] (a) edge (d2);
      \draw[tuh] (d2) edge (b);
      \node at (0,-2.5) {$G$};
    \end{scope}
    \begin{scope}[xshift=4cm,yshift=-8cm]
      \node[ndout] (a) at (0,0) {$a$};
      \node[ndout] (b) at (0,-1.5) {$b$};
      \node[ndout] (z) at (-1.5,0) {$z$};
      \node[ndint] (d1) at (-0.75,1) {$d_1$};
      \node[ndint] (d2) at (1,-0.75) {$d_2$};
      \draw[tuh] (a) edge (z);
      \draw[tuh] (d1) edge (z);
      \draw[huh] (a) edge (b);
      \draw[tuh] (d2) edge (b);
      \node at (0,-2.5) {$G_{\doit(X)}$};
    \end{scope}
    \begin{scope}[xshift=8cm,yshift=-8cm]
      \node[ndout] (a) at (0,0) {$a$};
      \node[ndout] (b) at (0,-1.5) {$b$};
      \node[ndout] (z) at (-1.5,0) {$z$};
      \node[ndint] (d1) at (-0.75,1) {$d_1$};
      \node[ndint] (d2) at (1,-0.75) {$d_2$};
      \draw[tuh] (a) edge (z);
      \draw[tuh] (d1) edge (z);
      \draw[huh] (a) edge (b);
      \draw[tuh] (d2) edge (b);
      \node at (0,-2.5) {$\MAG(G_{\doit(X)}) = H^*$};
    \end{scope}
    \begin{scope}[xshift=-4cm,yshift=-12cm]
      \node[ndout] (a) at (0,0) {$a$};
      \node[ndout] (b) at (0,-1.5) {$b$};
      \node[ndout] (z) at (-1.5,0) {$z$};
      \node[ndout] (d1) at (-0.75,1) {$d_1$};
      \node[ndout] (d2) at (1,-0.75) {$d_2$};
      \draw[tuh] (a) edge (z);
      \draw[tuh] (a) edge (d1);
      \draw[tuh] (d1) edge (z);
      \draw[tuh] (a) edge (b);
      \draw[tuh] (a) edge (d2);
      \draw[tuh] (d2) edge (b);
      \draw[huh] (z) edge (b);
      \node at (0,-2.5) {$H$};
    \end{scope}
    \begin{scope}[yshift=-12cm]
      \node[ndout] (a) at (0,0) {$a$};
      \node[ndout] (b) at (0,-1.5) {$b$};
      \node[ndout] (z) at (-1.5,0) {$z$};
      \node[ndout] (d1) at (-0.75,1) {$d_1$};
      \node[ndout] (d2) at (1,-0.75) {$d_2$};
      \draw[huh] (a) edge (z);
      \draw[tuh] (a) edge (d1);
      \draw[tuh] (d1) edge (z);
      \draw[huh] (a) edge (b);
      \draw[tuh] (a) edge (d2);
      \draw[tuh] (d2) edge (b);
      \node at (0,-2.5) {$G$};
    \end{scope}
    \begin{scope}[xshift=4cm,yshift=-12cm]
      \node[ndout] (a) at (0,0) {$a$};
      \node[ndout] (b) at (0,-1.5) {$b$};
      \node[ndout] (z) at (-1.5,0) {$z$};
      \node[ndint] (d1) at (-0.75,1) {$d_1$};
      \node[ndint] (d2) at (1,-0.75) {$d_2$};
      \draw[huh] (a) edge (z);
      \draw[tuh] (d1) edge (z);
      \draw[huh] (a) edge (b);
      \draw[tuh] (d2) edge (b);
      \node at (0,-2.5) {$G_{\doit(X)}$};
    \end{scope}
    \begin{scope}[xshift=8cm,yshift=-12cm]
      \node[ndout] (a) at (0,0) {$a$};
      \node[ndout] (b) at (0,-1.5) {$b$};
      \node[ndout] (z) at (-1.5,0) {$z$};
      \node[ndint] (d1) at (-0.75,1) {$d_1$};
      \node[ndint] (d2) at (1,-0.75) {$d_2$};
      \draw[huh] (a) edge (z);
      \draw[tuh] (d1) edge (z);
      \draw[huh] (a) edge (b);
      \draw[tuh] (d2) edge (b);
      \node at (0,-2.5) {$\MAG(G_{\doit(X)}) = H^*$};
    \end{scope}
  \end{tikzpicture}
  \caption{Counterexample to Zhang's claim: ``Furthermore, making this change will not make any other such directed edge in $\C{M}_{\underline{Y}\overline{X}}$ visible.''?
  $X = \{d_1,d_2\}$. 
  Well, it is not a counterexample to Zhang's claim, because his ``such directed edge'' is different (indeed, $G$-specific) from the way I interpreted it ($H$-specific).
  But it nicely shows that even if we can replace a single indirect and possibly-indirect directed edge (either $a \tuh z$ or $a \tuh b$) and maintain the same MAG representation ($H$),
  we cannot do the two replacements simultaneously, because that introduces additional induced paths / adjacencies.}
\end{figure}

What $H^*$ do exist?
\begin{Prp}
  Let $H$ be a DMAG with input nodes $J$ and output nodes $V$, that represents some CADMG with input nodes $J$ and output nodes $V^+ = V \dcup S$ given $S$.
  Let $D \subseteq V$.
  Define 
  $$\C{H}_{\doit(D)} := \{ \MAG(G_{\doit(D)} \given S) : \text{$G$ is a CADMG with input nodes $J$ and output nodes $V^+ = V \dcup S$ that is represented by $H$ given $S$} \}.$$
  Then:
  \begin{enumerate}
    \item there exists $H^* \in \C{H}_{\doit(D)}$ with $\skel(H^*) = \skel(H_{\bar{D}})$
    \item if $i \tuh j$ in $H$ is possibly-indirect, and there are no directed paths from $i$ to $j$ in $H_{\bar D}$ except for $i \tuh j$, there exists $H^* \in \C{H}_{\doit(D)}$ with $i \huh j$ in $H^*$.
  \end{enumerate}
\end{Prp}
\begin{proof}
  \begin{enumerate}
    \item If $H$ is a DMAG then $H$ is also a CADMG that is represented by $H$.
    \item Note: any edge $i \tuh j$ in $H$ \emph{could} be present in $G$ in the sense that there exists a $G$ represented by $H$ in which it is present, in which case it will end up in $H^*$ if $j \notin D$.
      If $i \tuh j$ in $H$, then note that if we take $G = H$ then we \emph{will} get $i \tuh j$ in $H^*$ if $j \notin D$.
    \item
      Indeed, if $i \tuh j$ it is possibly-indirect in $H$, there exists $G'$ represented by $H$ which doesn't contain $i \tuh j$.
      If $i \in \Anc_{{G'}_{\doit(D)}}(j)$, then there is a directed path from $i$ to $j$ in ${G'}_{\doit(D)}$,
      hence there is a directed path from $i$ to $j$ in $G'$ that avoids edges into $D$ (which is not $i \tuh j$),
      hence apart from $i \tuh j$ there is a directed path from $i$ to $j$ in $H$ that avoids edges into $D$, 
      hence apart from $i \tuh j$ there is a directed path from $i$ to $j$ in $H_{\bar D}$.
      Therefore, if there are no such paths, then $i \notin \Anc_{{G'}_{\doit(D)}}(j)$, and the corresponding $H^*$ has $i \huh j$.
%      Any directed path from $i$ to $j$ in $G'$ must also be present in $H$.
%      Any directed path from $i$ to $j$ in $G'$ that avoids edges into $D$ must also be present in $H_{\bar D}$.
%      So if there are no directed paths from $i$ to $j$ in $H_{\bar D}$ except for $i \tuh j$, then no directed paths from $i$ to $j$ that avoid edges into $D$ can be present in $G'$.
%      In that case, no directed paths from $i$ to $j$ can be present in ${G'}_{\doit(D)}$.
%      Hence $i \notin \Anc_{{G'}_{\doit(D)}}(j)$.
      \Joris{Note that if $i \tuh j$ is also invisible, as $G'$ we can take $H$ but with $i \tuh j$ replaced by $i \huh j$.
      It is clear that this $G'$ has an indirect directed path from $i$ to $j$.}
  \end{enumerate}
\end{proof}
We give the following tentative definition:
\begin{Def}\label{def:mag_do}
  Let $H$ be a MAG with input nodes $J$ and output nodes $V$, that represents some CADMG with input nodes $J$ and output nodes $V^+ = V \dcup S$ given $S$.
  Let $D \subseteq V$.
  Define $H_{\doit(D)}$ as the ``PAG'' with input nodes $J$ and output nodes $V$, such that 
  \begin{enumerate}
    \item the skeleton of $H_{\doit(D)}$ is that of $H_{\bar D}$
    \item every arrowhead in $H_{\bar D}$ is copied over to $H_{\doit(D)}$
    \item every tail on a definitely direct edge $i \tuh j$ in $H_{\bar D}$ is copied over to $H_{\doit(D)}$ (if $j \notin D$)
    \item for a possibly-indirect edge $i \tuh j$ in $H_{\bar D}$, if there are no directed paths from $i$ to $j$ in $H_{\bar D}$ except for $i \tuh j$, we put $i \ouh j$ in $H_{\doit(D)}$
    \item for a possibly-indirect edge $i \tuh j$ in $H_{\bar D}$, if there are directed paths from $i$ to $j$ in $H_{\bar D}$ except for $i \tuh j$, we put $i \tuh j$ in $H_{\doit(D)}$
  \end{enumerate}
\end{Def}
\begin{Prp}\label{prp:camdg_do_mag}
  Let $G$ be a CADMG with input nodes $J$ and output nodes $V^+ = V \dcup S$. 
  Denote the MAG that represents $G$ given $S$ by $H := \MAG(G \given S)$.
  For $D \subseteq V$, define $G^* := G_{\doit(D)}$ and $H^* := \MAG(G_{\doit(D)} \given S)$.
  Then $H^*$ is weakly represented by $H_{\doit(D)}$.
\end{Prp}
\begin{proof}
  Note that $x \in \Anc_{G_{\doit(D)}}(U)$ implies that $x \in \Anc_G(U)$, for all $x \in V^+ \cup J$ and $U \subseteq V^+ \cup J$.

  NOTE: any directed edge (path) in $G$ ends up in $M$ as well.
  So if in $M$ there is no directed path from $i$ to $j$ that avoids $D$ then in $G$ there is no such path either.
  We then know that in all possible $G_{\doit(D)}$, $i$ is not ancestor of $j$, and hence $i \huh j$ in $H^*$ if the two are adjacent in $H^*$.
  
  If $i \tuh j$ in $H$ is possibly-indirect and visible, then...we can recurse.
  Indeed, there must exist $k$ with $i \tuh k \tuh j$ in $M$ and $i \tuh k$ visible and $k \tuh j$ invisible in $M$.
  (We can pick $k$ to be a child of $i$ in $G$; every directed path from $i$ to $j$ in $G$ must go via a child; all children of $i$ in $G$ end up as such in $H$; so every possible directed path from $i$ to $j$ in $G$ must go via a child of $i$ in $H$)
  If both are definitely direct then they both end up in $H_{\doit(D)}$ and therefore also $i \tuh j$ must be in $H_{\doit(D)}$.
  TODO: Finish this argument!
  First note that every directed edge: if it is possibly-indirect we can replace it with a directed path, and otherwise it can be oriented as a directed edge in $H_{\doit(D)}$.
  For any already oriented directed path we can use transitivity to orient the ``shortcuts''.

  HM, the algorithm could become:
  \begin{enumerate}
    \item orient each definitely direct edge $i \tuh j$ (as $i \tuh j$);
    \item take the transitive closure;
  \end{enumerate}
  Question: if we can get $i \huh j$ in $H^*$, does this imply additional orientations?
  
  The case that remains is:
  if $i \tuh j$ in $H$ is possibly-indirect and visible, there must exist $k$ with $i \tuh k \tuh j$ in $H$ and $i \tuh k$ visible and $k \tuh j$ invisible in $H$.
  Three cases: all such $k$ are in $D$, some such $k$ are in $D$ and others are not in $D$, all such $k$ are not in $D$.
  We can further distinguish with respect to $i \tuh k$ definitely direct and $k \tuh j$ definitely direct.
  Note that we already know how to orient all the $k \tuh j$.
  If $i \tuh k$ is definitely direct we can also orient it.
  This means we will have $i \tuh k \tuh j$ in $H_{\doit(D)}$ (which implies $i \tuh j$ in $H_{\doit(D)}$), or $i \tuh k \ouh j$ in $H_{\doit(D)}$.
  Does the latter imply anything about the orientation of $i \suh j$ in $H_{\doit(D)}$?
  We know that $k \tuh j$ is possibly-indirect and invisible, so we can take for $G'$ the ADMG $H$ minus $i \tuh j$ plus $i \huh j$;
  could we have $i$ ancestor of $j$ in that $G'$? 
  If so, there must be $i \tuh k \huh j$ and a directed path $i \tuh \dots \tuh j$.
  This seems to suggest that there is also a $i \tuh k' \tuh j$ in $H$, with $k' \ne k$.
  We can similarly orient this. 
  Suppose we keep going, and end up with $k_1, \dots, k_n$ with for all $k_r$: $i \tuh k_r \tuh j$ in $H$ with $i \tuh k_r$ visible and $k_r \tuh j$ invisible,
  and we have oriented $i \tuh k_r \ouh j$ in $H_{\doit(D)}$ for all $k_r$.

  Then we have the final remaining interesting case.
  Perhaps this is a contradiction?
  Because shouldn't there at least remaing ONE ``non-spurious'' ancestral path?
  Ah no, the point is that this might get blocked.
  Perhaps we need to check now whether all $k_r$ are in $D$, or only some, or none.
  If none are in $D$, then we know that in each $G$ at least one directed path from $i$ to $j$ exists that won't disappear by $\doit(D)$.
  If all are in $D$, then I think we can conclude that after $\doit(D)$ $i$ is no longer ancestor of $j$.
  Finally, if some of them are in $D$, then we don't know?
  Wait, if there's at least one $k_r \in D$ then we can construct an ADMG in which the only directed path from $i$ to $j$ goes via $k_r$, which disappears under $\doit(D)$, and therefore we have to orient $i \ouh j$ in $H_{\doit(D)}$.

  So it looks like we're done, if we can show some lemmata to hold
  (that we can actually pick the ADMG in the last sentence of the previous paragraph, and that all directed paths from $i$ to $j$ have to go via one of the $k_r$).
\end{proof}

\Joris{
The next thing is to show (e.g.\ via Zhang's approach) that the circles don't matter for the m-separations.
Alternatively, we could just make use of the definite m-connections/separations.
Zhang essentially shows that we can treat the circles as tails for the purposes of separations.

Zhang (2008, Lemma 12) shows that (without selection bias) a directed edge in $H^*$ must be present in $H_{\bar D}$, but
that a bidirected edge in $H^*$ is either a bidirected edge in $H_{\bar D}$ or an invisible directed edge in $H_{\bar D}$.
Then he shows that $H_{\bar D}$ can be transformed into a supergraph of $H^*$ via a sequence of equivalence-preserving mark changes (Zhang and Spirtes 2005, Tian 2005).
Apparently we need Lemma 1 in Zhang \& Spirtes (2005).
This says:
For a DMAG $M$ and directed edge $a \tuh b$ in $M$, let $M'$ be the graph obtained from $M$ by replacing $a \tuh b$ by $a \huh b$.
Then $M'$ is a DMAG and Markov equivalent to $M$ if:
\begin{enumerate}
   \item there is no directed path in $G$ from $a$ to $b$ other than $a \tuh b$;
   \item for any $c \tuh a$ in $M$, $c \tuh b$ is also in $M$;
     for any $d \huh a$ in $M$, there is $d \suh b$ in $M$.
   \item there is no discriminating path for $a$ on which $b$ is the endpoint adjacent to $a$ in $M$.
\end{enumerate}
The proof mostly establishes that $M'$ is a valid DMAG; once that has been shown, the characterisation for Markov equivalence of DMAGs can be used.}

Is it true that any ancestral relation present in a MAG implies the ancestral relation is present in all CADMGs it represents?
Not necessarily!
What is true is that \emph{anterior} relations in a MAG correspond with ancestral relations in (all) CADMGs it represents.

\begin{figure}\centering
  \begin{tikzpicture}
    \begin{scope}
      \node at (0.0,0.75) {$G^1$:};
      \node[ndout] (X1) at (-1.5,0) {$i$};
      \node[ndout] (X2) at (0,0) {$d$};
      \node[ndout] (X3) at (1.5,0) {$j$};
%      \node[ndout] (X4) at (-3,0) {$l$};
%      \draw[tuh] (X4) edge (X1);
      \draw[tuh] (X1) edge (X2);
      \draw[tuh] (X2) edge (X3);
      \draw[huh, bend right] (X1) edge (X3);
    \end{scope}
    \begin{scope}[xshift=6cm]
      \node at (0,0.75) {$G^1_{\doit(d)}$: \Joris{($G^*$)}};
      \node[ndout] (X1) at (-1.5,0) {$i$};
      \node[ndint] (X2) at (0,0) {$d$};
      \node[ndout] (X3) at (1.5,0) {$j$};
%      \node[ndout] (X4) at (-3,0) {$l$};
%      \draw[tuh] (X4) edge (X1);
      \draw[tuh] (X2) edge (X3);
      \draw[huh, bend right] (X1) edge (X3);
    \end{scope}
    \begin{scope}[xshift=12cm]
      \node at (0,0.75) {$\MAG(G^1_{\doit(d)})$: \Joris{($H^*$)}};
      \node[ndout] (X1) at (-1.5,0) {$i$};
      \node[ndint] (X2) at (0,0) {$d$};
      \node[ndout] (X3) at (1.5,0) {$j$};
%      \node[ndout] (X4) at (-3,0) {$l$};
%      \draw[tuh] (X4) edge (X1);
      \draw[tuh] (X2) edge (X3);
      \draw[huh, bend right] (X1) edge (X3);
    \end{scope}
    \begin{scope}[yshift=-2cm]
      \node at (0.0,0.75) {$G^2$:};
      \node[ndout] (X1) at (-1.5,0) {$i$};
      \node[ndout] (X2) at (0,0) {$d$};
      \node[ndout] (X3) at (1.5,0) {$j$};
%      \node[ndout] (X4) at (-3,0) {$l$};
%      \draw[tuh] (X4) edge (X1);
      \draw[tuh] (X1) edge (X2);
      \draw[tuh] (X2) edge (X3);
      \draw[tuh, bend right] (X1) edge (X3);
    \end{scope}
    \begin{scope}[xshift=6cm,yshift=-2cm]
      \node at (0,0.75) {$G^2_{\doit(d)}$:};
      \node[ndout] (X1) at (-1.5,0) {$i$};
      \node[ndint] (X2) at (0,0) {$d$};
      \node[ndout] (X3) at (1.5,0) {$j$};
%      \node[ndout] (X4) at (-3,0) {$l$};
%      \draw[tuh] (X4) edge (X1);
      \draw[tuh] (X2) edge (X3);
      \draw[tuh, bend right] (X1) edge (X3);
    \end{scope}
    \begin{scope}[xshift=12cm,yshift=-2cm]
      \node at (0,0.75) {$\MAG(G^2_{\doit(d)})$:};
      \node[ndout] (X1) at (-1.5,0) {$i$};
      \node[ndint] (X2) at (0,0) {$d$};
      \node[ndout] (X3) at (1.5,0) {$j$};
%      \node[ndout] (X4) at (-3,0) {$l$};
%      \draw[tuh] (X4) edge (X1);
      \draw[tuh] (X2) edge (X3);
      \draw[tuh, bend right] (X1) edge (X3);
    \end{scope}
    \begin{scope}[yshift=-4cm]
      \node at (0.0,0.75) {$G^3$:};
      \node[ndout] (X1) at (-1.5,0) {$i$};
      \node[ndout] (X2) at (0,0) {$d$};
      \node[ndout] (X3) at (1.5,0) {$j$};
%      \node[ndout] (X4) at (-3,0) {$l$};
%      \draw[tuh] (X4) edge (X1);
      \draw[tuh] (X1) edge (X2);
      \draw[tuh] (X2) edge (X3);
      \draw[tuh, bend right] (X1) edge (X3);
      \draw[huh, bend left] (X1) edge (X3);
    \end{scope}
    \begin{scope}[xshift=6cm,yshift=-4cm]
      \node at (0,0.75) {$G^3_{\doit(d)}$:};
      \node[ndout] (X1) at (-1.5,0) {$i$};
      \node[ndint] (X2) at (0,0) {$d$};
      \node[ndout] (X3) at (1.5,0) {$j$};
%      \node[ndout] (X4) at (-3,0) {$l$};
%      \draw[tuh] (X4) edge (X1);
      \draw[tuh] (X2) edge (X3);
      \draw[tuh, bend right] (X1) edge (X3);
      \draw[huh, bend left] (X1) edge (X3);
    \end{scope}
    \begin{scope}[xshift=12cm,yshift=-4cm]
      \node at (0,0.75) {$\MAG(G^3_{\doit(d)})$:};
      \node[ndout] (X1) at (-1.5,0) {$i$};
      \node[ndint] (X2) at (0,0) {$d$};
      \node[ndout] (X3) at (1.5,0) {$j$};
%      \node[ndout] (X4) at (-3,0) {$l$};
%      \draw[tuh] (X4) edge (X1);
      \draw[tuh] (X2) edge (X3);
      \draw[tuh, bend right] (X1) edge (X3);
    \end{scope}
    \begin{scope}[yshift=-6.5cm]
      \node at (0,0.75) {MAG $H$:};
      \node[ndout] (X1) at (-1.5,0) {$i$};
      \node[ndout] (X2) at (0,0) {$d$};
      \node[ndout] (X3) at (1.5,0) {$j$};
%      \node[ndout] (X4) at (-3,0) {$l$};
%      \draw[tuh] (X4) edge (X1);
      \draw[tuh] (X1) edge (X2);
      \draw[tuh] (X2) edge (X3);
      \draw[tuh, bend right] (X1) edge (X3);
    \end{scope}
    \begin{scope}[yshift=-6.5cm,xshift=6cm]
      \node at (0,0.75) {MAG $H_{\bar d}$:};
      \node[ndout] (X1) at (-1.5,0) {$i$};
      \node[ndint] (X2) at (0,0) {$d$};
      \node[ndout] (X3) at (1.5,0) {$j$};
%      \node[ndout] (X4) at (-3,0) {$l$};
%      \draw[tuh] (X4) edge (X1);
      \draw[tuh] (X2) edge (X3);
      \draw[tuh, bend right] (X1) edge (X3);
    \end{scope}
    \begin{scope}[yshift=-6.5cm,xshift=12cm]
      \node at (0,0.75) {``PAG $H_{\doit(d)}$'':};
      \node[ndout] (X1) at (-1.5,0) {$i$};
      \node[ndint] (X2) at (0,0) {$d$};
      \node[ndout] (X3) at (1.5,0) {$j$};
%      \node[ndout] (X4) at (-3,0) {$l$};
%      \draw[tuh] (X4) edge (X1);
      \draw[tuh] (X2) edge (X3);
      \draw[ouh, bend right] (X1) edge (X3);
    \end{scope}
  \end{tikzpicture}
  \caption{MAG $H$ represents three different CADMGs $G^1, G^2, G^3$.
    After a hard intervention $\doit(d)$, their representative MAGs no longer coincide.
    MAG $H_{\bar d}$ represents $\MAG(G^2_{\doit(d)})$ but does not represent $\MAG(G^1_{\doit(d)})$.
    Note: there is no definite $m$-connecting path between $l$ and $j$ in the ``PAG $H_{\doit(d)}$'' 
    (which would lead us to the wrong conclusion that $l$ and $j$ are $id$-separated in each $G_{\doit(D)}$),
    but there is a definite $m$-connecting path between $l$ and $j$ in the MAG $H_{\bar d}$.
    \label{fig:mag_do}}
\end{figure}

\begin{figure}\centering
  \begin{tikzpicture}
    \begin{scope}
      \node at (0.0,0.75) {$G^1$:};
      \node[ndout] (X1) at (-1.5,0) {$i$};
      \node[ndout] (X2) at (0,0) {$d$};
      \node[ndout] (X3) at (1.5,0) {$j$};
      \node[ndout] (X4) at (-3,0) {$l$};
      \draw[tuh] (X4) edge (X1);
      \draw[tuh] (X1) edge (X2);
      \draw[tuh] (X2) edge (X3);
      \draw[huh, bend right] (X1) edge (X3);
    \end{scope}
    \begin{scope}[xshift=6cm]
      \node at (0,0.75) {$G^1_{\doit(d)}$: \Joris{($G^*$)}};
      \node[ndout] (X1) at (-1.5,0) {$i$};
      \node[ndint] (X2) at (0,0) {$d$};
      \node[ndout] (X3) at (1.5,0) {$j$};
      \node[ndout] (X4) at (-3,0) {$l$};
      \draw[tuh] (X4) edge (X1);
      \draw[tuh] (X2) edge (X3);
      \draw[huh, bend right] (X1) edge (X3);
    \end{scope}
    \begin{scope}[xshift=12cm]
      \node at (0,0.75) {$\MAG(G^1_{\doit(d)})$: \Joris{($H^*$)}};
      \node[ndout] (X1) at (-1.5,0) {$i$};
      \node[ndint] (X2) at (0,0) {$d$};
      \node[ndout] (X3) at (1.5,0) {$j$};
      \node[ndout] (X4) at (-3,0) {$l$};
      \draw[tuh] (X4) edge (X1);
      \draw[tuh] (X2) edge (X3);
      \draw[huh, bend right] (X1) edge (X3);
    \end{scope}
    \begin{scope}[yshift=-2cm]
      \node at (0.0,0.75) {$G^2$:};
      \node[ndout] (X1) at (-1.5,0) {$i$};
      \node[ndout] (X2) at (0,0) {$d$};
      \node[ndout] (X3) at (1.5,0) {$j$};
      \node[ndout] (X4) at (-3,0) {$l$};
      \draw[tuh] (X4) edge (X1);
      \draw[tuh] (X1) edge (X2);
      \draw[tuh] (X2) edge (X3);
      \draw[tuh, bend right] (X1) edge (X3);
      \draw[tuh, bend left] (X4) edge (X3);
    \end{scope}
    \begin{scope}[xshift=6cm,yshift=-2cm]
      \node at (0,0.75) {$G^2_{\doit(d)}$:};
      \node[ndout] (X1) at (-1.5,0) {$i$};
      \node[ndint] (X2) at (0,0) {$d$};
      \node[ndout] (X3) at (1.5,0) {$j$};
      \node[ndout] (X4) at (-3,0) {$l$};
      \draw[tuh] (X4) edge (X1);
      \draw[tuh] (X2) edge (X3);
      \draw[tuh, bend right] (X1) edge (X3);
      \draw[tuh, bend left] (X4) edge (X3);
    \end{scope}
    \begin{scope}[xshift=12cm,yshift=-2cm]
      \node at (0,0.75) {$\MAG(G^2_{\doit(d)})$:};
      \node[ndout] (X1) at (-1.5,0) {$i$};
      \node[ndint] (X2) at (0,0) {$d$};
      \node[ndout] (X3) at (1.5,0) {$j$};
      \node[ndout] (X4) at (-3,0) {$l$};
      \draw[tuh] (X4) edge (X1);
      \draw[tuh] (X2) edge (X3);
      \draw[tuh, bend right] (X1) edge (X3);
      \draw[tuh, bend left] (X4) edge (X3);
    \end{scope}
    \begin{scope}[yshift=-4cm]
      \node at (0.0,0.75) {$G^3$:};
      \node[ndout] (X1) at (-1.5,0) {$i$};
      \node[ndout] (X2) at (0,0) {$d$};
      \node[ndout] (X3) at (1.5,0) {$j$};
      \node[ndout] (X4) at (-3,0) {$l$};
      \draw[tuh] (X4) edge (X1);
      \draw[tuh] (X1) edge (X2);
      \draw[tuh] (X2) edge (X3);
      \draw[huh, bend right] (X1) edge (X3);
      \draw[tuh, bend left] (X4) edge (X3);
    \end{scope}
    \begin{scope}[xshift=6cm,yshift=-4cm]
      \node at (0,0.75) {$G^3_{\doit(d)}$:};
      \node[ndout] (X1) at (-1.5,0) {$i$};
      \node[ndint] (X2) at (0,0) {$d$};
      \node[ndout] (X3) at (1.5,0) {$j$};
      \node[ndout] (X4) at (-3,0) {$l$};
      \draw[tuh] (X4) edge (X1);
      \draw[tuh] (X2) edge (X3);
      \draw[huh, bend right] (X1) edge (X3);
      \draw[tuh, bend left] (X4) edge (X3);
    \end{scope}
    \begin{scope}[xshift=12cm,yshift=-4cm]
      \node at (0,0.75) {$\MAG(G^3_{\doit(d)})$:};
      \node[ndout] (X1) at (-1.5,0) {$i$};
      \node[ndint] (X2) at (0,0) {$d$};
      \node[ndout] (X3) at (1.5,0) {$j$};
      \node[ndout] (X4) at (-3,0) {$l$};
      \draw[tuh] (X4) edge (X1);
      \draw[tuh] (X2) edge (X3);
      \draw[huh, bend right] (X1) edge (X3);
      \draw[tuh, bend left] (X4) edge (X3);
    \end{scope}
    \begin{scope}[yshift=-6.5cm]
      \node at (0,0.75) {MAG $H$:};
      \node[ndout] (X1) at (-1.5,0) {$i$};
      \node[ndout] (X2) at (0,0) {$d$};
      \node[ndout] (X3) at (1.5,0) {$j$};
      \node[ndout] (X4) at (-3,0) {$l$};
      \draw[tuh] (X4) edge (X1);
      \draw[tuh] (X1) edge (X2);
      \draw[tuh] (X2) edge (X3);
      \draw[tuh, bend right] (X1) edge (X3);
      \draw[tuh, bend left] (X4) edge (X3);
    \end{scope}
    \begin{scope}[yshift=-6.5cm,xshift=6cm]
      \node at (0,0.75) {MAG $H_{\bar d}$:};
      \node[ndout] (X1) at (-1.5,0) {$i$};
      \node[ndint] (X2) at (0,0) {$d$};
      \node[ndout] (X3) at (1.5,0) {$j$};
      \node[ndout] (X4) at (-3,0) {$l$};
      \draw[tuh] (X4) edge (X1);
      \draw[tuh] (X2) edge (X3);
      \draw[tuh, bend right] (X1) edge (X3);
      \draw[tuh, bend left] (X4) edge (X3);
    \end{scope}
    \begin{scope}[yshift=-6.5cm,xshift=12cm]
      \node at (0,0.75) {``PAG $H_{\doit(d)}$'':};
      \node[ndout] (X1) at (-1.5,0) {$i$};
      \node[ndint] (X2) at (0,0) {$d$};
      \node[ndout] (X3) at (1.5,0) {$j$};
      \node[ndout] (X4) at (-3,0) {$l$};
      \draw[tuh] (X4) edge (X1);
      \draw[tuh] (X2) edge (X3);
      \draw[ouh, bend right] (X1) edge (X3);
      \draw[tuh, bend left] (X4) edge (X3);
    \end{scope}
  \end{tikzpicture}
  \caption{MAG $H$ represents three different CADMGs $G^1, G^2, G^3$.
    After a hard intervention $\doit(d)$, their representative MAGs no longer coincide.
    \label{fig:mag_do2}}
\end{figure}

\subsubsection{First attempts at generalizing to (cyclic) $\sigma$-setting with selection bias}

\Joris{
Here is a list of open problems:
\begin{enumerate}
  \item A complete way to read off which nodes in a $\sigma$-MAG are non-ancestors of selection nodes.
    (Conjecture: all nodes that have an adjacent arrowhead, or for which there is an undirected path to a node with an adjacent arrowhead.)
  \item A complete way to read off which nodes in a $\sigma$-MAG are ancestors of selection nodes.
    (Partial solution: some FCI rules allow us to identify this.)
  \item With these building blocks, we may be able to extend the power of the intervention node formalism we now have for CDMGs to $\sigma$-MAGs.
\end{enumerate}
}

\Joris{
  \begin{Def}
    Let $G$ be a CDMG with input nodes $J$ and output nodes $V^+ = V \dcup S$ and $H$ a $\sigma$-MAG that represents $G$ given $S$.
    Let $a,b \in V \cup J$. FOR NOW, ASSUME $J = \emptyset$.
    An undirected edge $a \tut b$ in $H$ is called \emph{audible} if there exists a path $a \tut b \tut v_1 \dots \tut v_n \hus c$ in $H$ (with possibly $n = 0$).
    Otherwise, $a \tut b$ is called \emph{inaudible}.
    A node $a \in V \cup J$ is called \emph{audible} if it is part of an audible edge, and \emph{inaudible} otherwise.
  \end{Def}
  Note that we cannot have an inaudible undirected edge adjacent to an audible undirected edge.
}

\begin{Con}
  Let $G$ be a CDMG with input nodes $J$ and output nodes $V^+ = V \dcup S$ and $H$ a $\sigma$-MAG that represents $G$ given $S$.
  Let $a \in V \cup J$. \Joris{For now: assume $J = \emptyset$!}
  There exists a path $a \tut \dots \tut b \hus c$ in $H$ (with possibly $a = b$) (Let us call this property ``$a$ is audible''!) implies $a \notin \Anc_G(S)$.
  If there doesn't exist such a path, then $a$ is possibly in $\Anc_G(S)$, that is, there exist $(G_1,S_1)$ and $(G_2,S_2)$ that are both represented by $H$ but where $a \in \Anc_{G_1}(S_1)$ but $a \notin \Anc_{G_2}(S_2)$.
\end{Con}
\begin{proof}
  $\implies$ is OK.

  Take a single node $a$. There doesn't exist such a path, and indeed, either $a$ is ancestor of $S$ or not.
  Can we find a CDMG that is represented by $H$ for which $a \notin \Anc(S)$?
  That is precisely what we were trying to do below in Lemma 12.10.3!

  If we manage to do so, then we can just add selection on $a$ to that CDMG to obtain one where $a \in \Anc(S)$.
\end{proof}

\Joris{
  Wait. Is the notion of inducing path in a $\sigma$-MAG correct? 
  In a $\sigma$-PAG it is OK (I believe), as there we can never conclude the presence of a cycle.
  But in a $\sigma$-MAG we can consider audible nodes / edges.
  A sequence of audible edges should also be considered inducing in a $\sigma$-MAG!
  Well, perhaps not, since then there will also be an edge between the two nodes in the $\sigma$-MAG.

  Also, we should be more liberal to colliders being ancestor of something: the directed path may also contain audible edges.

  So we should consider more paths in a $\sigma$-MAG as inducing!
  This doesn't seem to invalidate any of the earlier results, as the only thing we've used so far is that if there is a definitely inducing path in a $\sigma$-PAG between two nodes, they must be adjacent.
  But we can apparently strengthen this result for $\sigma$-MAGs (at least).
  It even appears that if we do so, we can show that we can find a ``canonical'' representative CDMG and set of selection nodes for a $\sigma$-MAG.
  This then in turn may allow us to read off conditional independences from a $\sigma$-MAG.
  (Assuming completeness of extended FCI, this should work, as all CDMGs represented by a $\sigma$-MAG should be Markov equivalent, as they are part of the FCI $\sigma$-PAG.)

  So let's try to define this. 
  Let $H$ be a $\sigma$-MAG that represents a CDMG $G$ given $S$.
  Call $v$ an \emph{ancestor} of $w$ in $H$ ($v \in \Anc_H(w)$) if there exists a path of audible edges and directed edges pointing towards $w$.
  Proposal: a walk between $v,w \in V \cup J$ (with $v \ne w$) in a $\sigma$-MAG $H$ (representing a CDMG $G$ given $S$) is called \emph{$\sigma$-inducing} if every non-endpoint node is a collider in $\Anc_H(\{v,w\})$.
  Hm, not sure if that's correct. Consider this: $a \suh b \tut c$ (with $b \tut c$ audible) should be inducing.
  Also $a \tut b \tut c$ (with both $a \tut b$, $b \tut c$ audible) should be inducing.
  Also $a \hut b \tut c$ (with $b \tut c$ audible) should be inducing.
  So basically every audible edge is inducing.
  Well, sometimes we can also just look for another inducing path, which would exist in all these cases (a single edge)!
  So it's not clear whether we need to adapt this definition.

  I see two options to proceed:
  \begin{itemize}
    \item Use the notion of $s$-separation and turn the audible undirected edges into ``equality'' edges?
    \item Keep using undirected edges for both selection bias and cycles, but adapt the notion of ``ancestor'' in a $\sigma$-MAG in such a way that ancestral relations are preserved;
      if this is possible, think about how to define inducing paths such that they are basically unblockable walks that can be concatenated together.
      Maximality should then come for free. The notion of ``inducing path'' in a $\sigma$-MAG may not be so important after all; what is more important is 
  \end{itemize}
}

\Joris{Here are some key definitions from above, for reference:
  A $\sigma$-PAG is defined by:
  \begin{enumerate}
    \item Between any two distinct nodes there is at most one edge, and there are no edges between a node and itself;
    \item No input node is adjacent to any other input node;
    \item The graph contains no directed cycles, no almost directed cycles, and no unshielded triple of the form $i \suh j \tut k$ (``$\sigma$-ancestral'');
    \item There is no definitely inducing path between any two distinct non-adjacent nodes (``maximal'');
%    \item If there is an edge $j \sus v$ with $j \in J, v \in V$ then it must be of the form $j \tus v$.
  \end{enumerate}
  $H$ represents $G$ given $S$ iff:
  \begin{enumerate}
    \item Between any two distinct nodes in $H$ there is at most one edge, and there are no edges between a node and itself;
    \item Two distinct nodes $i,j \in V \cup J$ are adjacent in $H$ if and only if $\{i,j\} \not\subseteq J$ and there is a $\sigma$-inducing path between $i$ and $j$ given $S$ in $G$;
    \item If $i \hus j$ in $H$ then $i \notin \Anc_{G}(\{j\} \cup S)$;
    \item If $i \tus j$ in $H$ then $i \in \Anc_{G}(\{j\} \cup S)$;
%    \item If $j \sus v$ in $H$ with $j \in J, v \in V$ then it must be of the form $j \tus v$.
  \end{enumerate}
  If $H$ represents $G$ given $S$, then $H$ is a $\sigma$-PAG.
  If $H$ has no circles, it is a $\sigma$-MAG.
}

The following notion was proposed originally by \cite{Zhang2008_JMLR}. We generalize it here to our setting (which allows for selection bias, input nodes and cycles).
\begin{Def}\label{def:mag_visible_edge_alt}
  Let $H$ be a $\sigma$-PAG with input nodes $J$ and output nodes $V$.

  Let $a \in J \cup V$ and $b \in V$.
  A directed edge $a \tuh b$ in $H$ is called \emph{visible} if $a \in J$, or if $a \in V$ and there is a node $z$ in $H$ that is not adjacent to $b$ such that:
  \begin{itemize}
    \item there is an edge $z \suh a$ in $H$, or 
    \item there is a definite collider path $z \suh \dots \huh a$ in $H$ that is into $a$ and every node on the path (except $z$) is a parent of $b$ in $H$. 
  \end{itemize}
  Otherwise, $a \tuh b$ is called \emph{invisible}.
\end{Def}

The following Lemma of \cite{Zhang2008_JMLR} stays valid in our more general setting.
\begin{Lem}\label{lem:zhang2008_lemma_9}
  Let $G$ be a CDMG with input nodes $J$ and output nodes $V^+ = V \dcup S$ and $H$ a $\sigma$-PAG that represents $G$ given $S$.
  Let $a \in V \cup J$, $b \in V$ such that $a \tuh b$ is a directed edge in $H$. 
  If $a \tuh b$ is visible, then there exists no $\sigma$-inducing walk between $a$ and $b$ given $S$ in $G$ that is into $a$.
  In particular, $a \huh b$ is not in $G$.
\end{Lem}
\begin{proof}
  If $a \in J$, then there are no edges into $a$ in $G$ by definition, which proves the claim.
  So assume $a \in V$. We proceed by contradiction. Assume that in $G$ there exists an $\sigma$-inducing walk between $a$ and $b$ given $S$ that is into $a$.

  Let $z$ be a node in $V \sm \{a,b\}$.
  Consider two cases:
  \begin{itemize}
    \item
      If $z \suh a$ is in $H$, then there is an $\sigma$-inducing walk between $z$ and $a$ in $G$ given $S$ that is into $a$ (Lemma~\ref{lem:mixed_graph_represents_CDMG}).
  Concatenating this with the $\sigma$-inducing walk between $a$ and $b$ given $S$ that is into $a$ we obtain an $\sigma$-inducing walk between $z$ and $b$ given $S$ (note that $a$ becomes a collider that is in $\Anc_G(\{b\} \cup S)$) in $G$. 
\item
  Suppose there is a definite collider path $\pi$ in $H$ from $z$ to $a$ that is into $a$ and such that every node on the path (except $z$) is parent of $b$ in $H$. 
  For each pair of subsequent nodes $(v_i,v_{i+1})$ on $\pi$ there exists an $\sigma$-inducing walk in $G$ between $v_i$ and $v_{i+1}$ given $S$ that is into $v_{i+1}$, and also into $v_i$ unless possibly $v_i = z$.
  Because each node other than $z$ on $\pi$ is in $\Anc_G(\{b\} \cup S)$, all these $\sigma$-inducing walks can be concatenated into one $\sigma$-inducing walk in $G$ between $z$ and $b$ given $S$.
  \end{itemize}
  \Joris{Note that we might be able to weaken the part ``$v_i$ is parent of $b$'' to ``$v_i \tus b$''? It seems that this can still be excluded in the cyclic case: then there would be
  an edge $z \suh v_1$ but no edge between $z$ and $b$ even though $z,b$ are in the same scc. So it seems parent is ok.}

  In both cases, $z$ and $b$ are adjacent in $H$.
  Since this holds for all such $z$, the directed edge $a \to b$ in $H$ cannot be visible. 
\end{proof}

%The following lemma shows that we can always interpret a $\sigma$-PAG as representing a CADMG given some latent selection nodes.
%NO THAT IS NOT POSSIBLE.
\Joris{The following works in the acyclic setting, but I haven't yet been able to prove this conjecture when allowing for cycles.
In the setting considered by \cite{Zhang2008_JMLR} (no input nodes, no cycles, no selection bias), the following statement is
trivial (I don't think Zhang bothered writing it down) because one can just take $\tilde{G} = H$.}
\begin{Lem}
  Let $G$ be a CDMG with input nodes $J$ and output nodes $V^+ = V \dcup S$ and $H$ the $\sigma$-MAG that represents $G$ given $S$.
  There exists a CDMG $\tilde{G}$ with input nodes $J$ and output nodes $V^+ = V \dcup \tilde{S}$ that is also represented by $H$ given $\tilde{S}$,
  and such that for each edge $i \tuh j$ in $H$, the edge $i \huh j$ is not in $\tilde{G}$.
  \Joris{Originally, I hoped that we can even take $\tilde{G}$ to be a CADMG, but that won't work: 
  If $G$ consists of an scc and a parent of that scc, then its $\sigma$-MAG will contain shielded triples $\tuh \tut$, something that we will never get from an acyclic $\tilde{G}$.
  It seems we first have to decide how we can identify/distinguish in a $\sigma$-MAG the ancestors of $S$, and the scc's.}
\end{Lem}
\begin{proof}
  Replace all undirected edges $x \tut y$ in $H$ with $x \tuh s_{\{x,y\}} \hut y$.
  Let $\tilde{S} := \{\{x,y\} : x \tut y \in H\}$.
  The result will be $\tilde{G}$.
  Note that $\tilde{G}$ is acyclic by construction.
  \Joris{This doesn't work if we allow for cycles. 
  Instead, try this: for undirected edges $a \tut b$ in $H$ which are part of a path $a \tut b \tut \dots \tut c \hus$, that is, a path of undirected edges with eventually an adjacent arrowhead, we know that all nodes on this path are non-ancestors of $S$ in $G$, and must be part of an scc (for which all its nodes are non-ancestors of $S$ in $G$).
  The nodes in this scc of $G$ must from an undirected clique in $H$.
  In $\tilde{G}$, we will turn this into a fully connected directed clique (with a directed edge between each ordered pair of nodes).
  All remaining undirected edges all undirected edges $x \tut y$ in $H$ will be replaced by $x \tuh s_{\{x,y\}} \hut y$.}

  Since all directed edges in $H$ are also present in $\tilde{G}$, and no directed edge absent in $H$ is present in $\tilde{G}$, the ancestral relations between variables in $V$ are the same in $H$ as in $\tilde{G}$. \Joris{This only holds in the acyclic setting; in the cyclic setting, $\tilde{G}$ contains all directed edges of $H$, and possibly more.}

  We will show that $H$ represents $\tilde{G}$ given $\tilde{S}$.
  For that we need to show:
  \begin{enumerate}
    \item For all $i \in V \cup J, j \in V$: $i \sus j$ is in $H$ if and only if there is a $\sigma$-inducing path between $i$ and $j$ given $\tilde{S}$ in $\tilde{G}$;
    \item If $i \hus j$ in $H$ then $i \notin \Anc_{\tilde{G}}(\{j\} \cup \tilde{S})$;
    \item If $i \tus j$ in $H$ then $i \in \Anc_{\tilde{G}}(\{j\} \cup \tilde{S})$;
  \end{enumerate}

  Note that for $v \in V$: if $v \notin \Anc_G(S)$ then $v \notin \Anc_{\tilde{G}}(\tilde{S})$.
  Or, equivalently: if $v \in \Anc_{\tilde{G}}(\tilde{S})$, then $v \in \Anc_G(S)$.
  (Note: the converse does not hold).
  Indeed, if $v \in \Anc_{\tilde{G}}(\tilde{S})$ then there must be a directed path $v \tuh \dots \tuh a \tuh s_{a,b}$ in $\tilde{G}$ such that all nodes except for $s_{a,b}$ are in $V$, and $s_{a,b} \in \tilde{S}$ for some edge $a \tut b \in H$.
  The directed path $v \tuh \dots \tuh a$ must be in $H$. 
  \Joris{But since $\suh a \tut b$ cannot occur in $H$ if $G$ is acyclic, the path must be the trivial path $v=a$.
  Hence $v \tut b$ in $H$.}
  %Hence $v \in \Anc_H(a)$, so $v \in \Anc_G(\{a\} \cup S)$.
  %Since $a \in \Anc_G(S)$ \Joris{that only holds for acyclic $G$}, $v \in \Anc_G(S)$.
  Hence $v \in \Anc_G(S)$ \Joris{holds only for acyclic $G$}.

  First adjacency.
  Let $x \in V \cup J$, $y \in V$.
  Suppose $x$ and $y$ are not adjacent in $H$.
  We must show that in $\tilde{G}$ there is no $\sigma$-inducing path between $x$ and $y$ given $\tilde{S}$.
  Suppose there is one.
  Pick one of minimal length.
  It must be a collider path in which every collider is a $\tilde{G}$-ancestor of $\{x,y\} \cup \tilde{S}$.
  \Joris{Cyclic: a path in which every collider is a $\tilde{G}$-ancestor of $\{x,y\} \cup \tilde{S}$ and every non-endpoint non-collider is unblockable in $\tilde{G}$.}

  Suppose the path contains a node of $\tilde{S}$.
  In that case, the path will look like $x \dots a \tuh s_{a,b} \hut b \dots y$ with $a \tut b$ in $H$.
  But since all non-endpoint non-colliders must be unblockable (which neither $a$ nor $b$ is, as they both point to $s_{a,b}$ which is in a different scc in $\tilde{G}$), we conclude that $a = x$ and $b = y$.
  Hence, the path is $x \tuh s_{x,y} \hut y$ with $x \tut y$ in $H$. Contradiction.
  We conclude that the path can only contain nodes in $V$.

  Suppose the path contains an audible edge, that is, an edge between two nodes $a,b$ such that $a \tut b \in H$ is audible. 
  Then that edge is either $a \tuh b$ or $a \hut b$, and furthermore, $a,b \notin \Anc_G(S)$  and $a,b \notin \Anc_{\tilde{G}}(\tilde{S})$ (??).
  \Joris{TODO!}

  Suppose each of the edges on the path is present in $H$.
  Then each collider $k$ on the path must be a non-ancestor of $S$ in $G$.
  Each collider on the path is $\tilde{G}$-ancestor of $\{x,y\} \cup \tilde{S}$.
  Hence there is a directed path from the collider $k$ to $x$ in $\tilde{G}$, or to $y$, or to $\tilde{S}$.
  If there is a directed path from the collider $k$ to $x$ or $y$ in $\tilde{G}$, then $k \in \Anc_H(\{x,y\})$.
  Otherwise, there is a directed path from the collider $k$ to $\tilde{S}$ in $\tilde{G}$.
  It must be of the form $k \tuh \dots \tuh a \tuh s_{ab}$ for some $a \tut b \in H$, so with $a \in \Anc_G(S)$.
  The directed path $k \tuh \dots \tuh a$ must also exist in $H$, and hence $k$ is a $G$-ancestor of $a$.
  Hence $k$ is $G$-ancestor of $S$. Contradiction.
  So then the path is also a definite inducing path between $x$ and $y$ in $H$, which contradicts the maximality of $H$.

  Hence, $x$ and $y$ are not adjacent in $MAG(\tilde{G})$ if $x$ and $y$ are not adjacent in $H$.

  Now suppose $x$ and $y$ are adjacent in $H$.

  Suppose $x \tuh y$ is in $H$.
  Then by construction $x \tuh y$ is in $\tilde{G}$.
  Hence it is in $MAG(\tilde{G})$.
  Then $y \notin \Anc_G(S)$, and therefore also $y \notin \Anc_{\tilde{G}}(\tilde{S})$.
  Hence $x \tuh y$ is in $MAG(\tilde{G})$.
  Similarly for the edge $x \hut y$ in $H$.

  If $x \huh y$ in $H$ then $x \huh y$ in $\tilde{G}$.
  Both $x, y \notin \Anc_G(S)$. 
  Hence, $x,y \notin \Anc_{\tilde{G}}(\tilde{S})$, and hence there is an edge $x \huh y$ in $MAG(\tilde{G})$.

  If $x \tut y$ in $H$ then $x \tuh s_{x,y} \hut y$ in $\tilde{G}$ and hence there is an edge $x \tut y$ in $MAG(\tilde{G})$.
\end{proof}

\begin{Con}
All this means that ``invisible'' is synonym for ``possibly-confounded'', and ``visible'' is synonym for ``unconfounded''.
This was known in the case: no cycles, no selection bias, no input nodes; we conjecture it remains true if we allow for cycles, selection bias and input nodes.
\end{Con}


\subsubsection{Other stuff}

The causal interpretation of the $\sigma$-PAG output by FCI is subtle, and not quite as transparent as that of a DMG.
When assuming no selection bias ($S = \emptyset$), the soundness of FCI allows us to directly read off some (non-)ancestral relations from the $\sigma$-PAG output by FCI. 
However, this is not yet all causal information that is identifiable from the $\sigma$-Markov equivalence class. 
With more effort, one can also read off additional (non-)ancestral relations, direct causal relations, the absence of confounding, and the absence of cycles. 

If selection bias can be ruled out, we can read off from the estimated $\sigma$-PAG consistent estimates for:
  \begin{enumerate}[label=(\roman*)]
%  \item the absence of causal cycles according to $\C{M}$ (via Proposition~\ref{prop:noncycles});
%  \item the absence/presence of (possibly indirect) causal relations according to $\C{M}$ (via Propositions~\ref{prop:cdpag_ancestors_cdmg} and \ref{prop:cdpag_nonancestors}); 
%  \item the absence of confounders according to $\C{M}$ (via Proposition~\ref{prop:identify_nonconfounders});
%  \item the absence/presence of direct causal relations according to $\C{M}$ (via Proposition~\ref{prop:identifiable_direct_causes}). 
%  \end{compactenum}
  \item the absence/presence of (possibly indirect) causal relations according to $M$; 
  \item the absence of latent common causes according to $M$;
  \item the absence/presence of direct causal relations according to $M$;
  \item the absence of causal cycles according to $M$,
  \end{enumerate}
as explained in detail in \cite{MooijClaassen_UAI_20}.
To the best of our knowledge, if selection bias is not ruled out, the causal interpretation of the $\sigma$-PAG output by FCI is not fully understood at present.
%Note that this corollary also applies to L-CBNs, as these can be considered to be special cases of simple SCMs.

Here we formulate some conjectures that have been proven under additional assumptions.

"\cite{Zhang2006} conjectured soundness and completeness of a criterion to read off all invariant ancestral relations from a complete DPAG.
Roumpelaki \emph{et al.} (2016) proved soundness of the criterion (and claim completeness, but their proof seems incomplete)."
\begin{Prp}[Proposition 4 in \cite{MooijClaassen_UAI_20}]
Let $G$ be a CDMG with input nodes $J$ and output nodes $V$, and $H$ be a $\sigma$-PAG that represents $G$.
Assume that all unshielded triples in $H$ have been oriented according to FCI rule $\C{R}0$ (and the background knowledge?) using $\IM_\sigma(G)$.
For two output nodes $i \ne j \in V$:
  \begin{enumerate}
    \item if there exists a directed path from $i$ to $j$ in $H$, or
    \item if there exist uncovered possibly directed paths from $i$ to $j$ in $H$ of the form $i, u, \dots, j$ and $i, v, \dots, j$ such that $u$ and $v$ are distinct non-adjacent output nodes in $H$
  \end{enumerate}
then $i \in \Anc_G(j)$.
\end{Prp}
\begin{proof}
  Mooij \& Claassen extended this result to DMGs.
  Still, no selection bias is assumed, and no input nodes either.
\end{proof}

\cite[p. 137]{Zhang2006} provides a sound and complete criterion to read off non-ancestors from a complete DPAG. 
\begin{Prp}[Proposition 5 in \cite{MooijClaassen_UAI_20}]
Let $G$ be a CDMG with input nodes $J$ and output nodes $V$, and $H$ be a $\sigma$-PAG that represents $G$.
For two output nodes $i \ne j \in V$:
if there is no possibly directed path from $i$ to $j$ in $H$, then $i \notin \Anc_G(j)$.
\end{Prp}
\begin{proof}
  Zhang proves this for ADMGs, Mooij \& Claassen extend it to DMGs.
  Two additional extensions are needed: selection bias, and input nodes.
\end{proof}
